\documentclass[10pt]{amsart}
\usepackage[a4paper,margin=23mm]{geometry}
\usepackage{amsmath,amssymb,amsthm,mathtools,hyperref,graphicx,xspace,letltxmacro,xpatch}
\usepackage[english]{babel}
\newtheorem{theo}{Theorem}[section]
\newtheorem{lemm}[theo]{Lemma}\newtheorem{prop}[theo]{Proposition}\newtheorem{coro}[theo]{Corollary}\newtheorem{conj}[theo]{Conjecture}
\theoremstyle{definition}\newtheorem{defi}[theo]{Definition}\newtheorem{rema}[theo]{Remark}\newtheorem{exam}[theo]{Example}\newtheorem{exem}[theo]{Example}\newtheorem{nota}[theo]{Notation}
\newcommand{\ie}{i.e.}\newcommand{\eg}{e.g.}\newcommand{\resp}{resp.}\newcommand{\skpt}{\par\smallskip\noindent}
\emergencystretch=3em\pagestyle{plain}

\usepackage[all,cmtip]{xy}

\DeclareMathOperator{\degkin}{deg_{\, kin}}
\DeclareMathOperator{\starkin}{\star_{\, kin}}




\newcommand{\N}{\mathbb{N}}
\newcommand{\R}{\mathbb{R}}
\newcommand{\T}{\mathbb{T}}
\newcommand{\bbS}{\mathbb{S}}
\newcommand{\bbP}{\mathbb{P}}


\newcommand{\cC}{\mathcal{C}}
\newcommand{\cL}{\mathcal{L}}
\newcommand{\cT}{\mathcal{T}}
\newcommand{\cG}{\mathcal{G}}
\newcommand{\cQ}{\mathcal{Q}}
\newcommand*\fR{\mathcal{R}}

\newcommand{\frka}{\mathfrak{a}}
\newcommand{\frkb}{\mathfrak{b}}
\newcommand{\frkc}{\mathfrak{c}}

\newcommand{\Qext}{\cQ_{\mathrm{ext}}}
\newcommand{\Qint}{\cQ_{\mathrm{int}}}


\DeclareMathOperator{\Id}{Id}


\newcommand*\eps{\varepsilon}
\newcommand*\un{\mathbf{1}}

\newcommand\dt{{\frac{\mathrm d}{\mathrm dt}}}
\newcommand{\dd}{{\, \mathrm d}}

%\newenvironment{step}[1]{\begin{proof}[#1]}{\let\qed\relax\end{proof}}
%\newenvironment{laststep}[1]{\begin{proof}[#1]}{\end{proof}}
%
%\newcounter{stepcount}
%%\renewcommand{\thestepcount}{\alph{stepcount}}
%\newenvironment{step}[1]{\refstepcounter{stepcount}\begin{proof}[Step~\thestepcount]}{\let\qed\relax\end{proof}}

\newcounter{stepcount}
\newenvironment{step}[1]{\refstepcounter{stepcount}\begin{proof}[Step~\thestepcount]{\emph{#1}}}{\let\qed\relax\end{proof}}

\newcounter{casecount}
\newenvironment{case}[1]{\refstepcounter{casecount}\begin{proof}[Case~\thecasecount]{\emph{#1}}}{\let\qed\relax\end{proof}}
%%%%%%%%%%%%%%%%%%%%%%%%%%%%%%%%%%%%%%%%%%%%%%%%%%%%%%%%%%%%%
\graphicspath{{./figures/}}

\newcommand*{\mk}{\mkern -1mu}
\newcommand*{\Mk}{\mkern -2mu}
\newcommand*{\mK}{\mkern 1mu}
\newcommand*{\MK}{\mkern 2mu}

\hypersetup{urlcolor=purple, linkcolor=blue, citecolor=red}

\newcommand*{\romanenumi}{\renewcommand*{\theenumi}{\roman{enumi}}}
\newcommand*{\Romanenumi}{\renewcommand*{\theenumi}{\Roman{enumi}}}
\newcommand*{\alphenumi}{\renewcommand*{\theenumi}{\alph{enumi}}}
\newcommand*{\Alphenumi}{\renewcommand*{\theenumi}{\Alph{enumi}}}
\let\oldtilde\tilde
\renewcommand*{\tilde}[1]{\mathchoice{\widetilde{#1}}{\widetilde{#1}}{\oldtilde{#1}}{\oldtilde{#1}}}
\let\oldforall\forall
\renewcommand*{\forall}{\mathrel{\oldforall}}
%%%%%%%%%%%%%%%%%%%%%%%%%%%%%%%%%%%%%%%%%%%%%%%%%%%%%%%%%%%%%

\title{Global well-posedness for a density-dependent kinetic Fokker-Planck model}
\author{Imbert, Cyril, Mouhot, Clément}
\date{doi:10.5802/ahl.75; CC BY 4.0}
\begin{document}
\maketitle
\noindent\textit{Editorial note.} Only selected author theorem counted; complete original mathematical body and bibliography preserved. Source-defined step/case environments retained; proofc class rendering supplied where needed.\par\medskip
\noindent\textit{Editorial note.} Publisher bibliography U+02BC transliteration apostrophe rendered with the standard apostrophe glyph.\par\medskip
\noindent\textit{Allowed standard prerequisites (editorial).}\begin{itemize}
\item Standard functional analysis, measure theory and the author-cited classical theorems. Standard prior results stay cited; complete source-local proofs and definitions retained.
\end{itemize}\medskip

\section{Introduction}

\subsection{The Schauder estimate for linear kinetic Fokker--Planck equations}

The first part of this paper deals with Schauder estimate for linear kinetic Fokker--Planck equations of the form
\begin{equation}\label{eq:linear}
\left(\partial_t + v \cdot \nabla_x\right) g =
\sum_{1\,\le\:i,\,j\:\le\,d} a^{i,\,j} \partial^2_{v_i\,v_j} g + \sum_{1\,\le\,i\,\le\,d}b^i \partial_{v_i} g + c g + S \quad \text{ in } \R \times \R^d \times \R^d
\end{equation}
for some given function $S$ under the assumption that the diffusion matrix\linebreak$A= (a^{i,\,j}(t,x,v))_{i,\,j=1,\,\dots,\,d}$ satisfy a \emph{uniform ellipticity condition} for some $\lambda>0$:
\begin{equation}\label{eq:elliptic}
\forall \, (t,x,v) \in \R \times \R^d \times \R^d, \
\forall \, \xi \in \R^d, \quad
\sum_{1\,\le\:i,\,j\:\le\,d} a^{i,\,j}(t,x,v) \xi_i \xi_j \ge \lambda |\xi|^2.
\end{equation}
The main result of this article is a \emph{Schauder estimate}, that is to say an a priori estimate for classical solutions to~\eqref{eq:linear} controlling their second-order H\"older regularity (in the sense of a ``kinetic order'' made precise below) by their supremum norm, under the assumptions that the coefficients $a^{i,j},b^i,c$ are also H\"older continuous.


\begin{theo}[The Schauder estimate]\label{thm:variable}
Given $\alpha \in (0,1)$ and $a^{i,\,j},b^i,c \in C^\alpha(\R \times \R^d \times \R^d)$, $i,j=1,\dots,d$, satisfying~\eqref{eq:elliptic} and a function $S \in C^\alpha(\R \times \R^d \times \R^d)$, any classical solution $g$ to~\eqref{eq:linear} satisfies
\begin{multline}\label{eq:schauder-main}
\left[\partial_t g + v \cdot \nabla_x g
\right]_{\cC^\alpha\left(\R\,\times\,\R^d\,\times\,\R^d\right)} + \left[\nabla_v ^2 g
\right]_{\cC^\alpha\left(\R\,\times\,\R^d\,\times\,\R^d\right)}\\
\le C \left([S]_{\cC^\alpha\left(\R\,\times\,\R^d\,\times\,\R^d\right)} +
\|g\|_{L^\infty\left(\R\,\times\,\R^d\,\times\,\R^d\right)}\right)
\end{multline}
where the constant $C$ depends on dimension $d$, the constant $\lambda$ from~\eqref{eq:elliptic}, the exponent $\alpha$ and the $[\cdot]_{\cC^\alpha}$ semi-norm and $L^\infty$ norm of $a^{i,\,j}$ for $i,j=1,\,\dots,\,d$, $b^i$ for $i=1,\,\dots,\,d$, and $c$. The semi-norm $[\cdot]_{\cC^\alpha}$ is the standard H\"older semi-norm for the distance $\|(t,x,v)\|=|t|^{\frac12} + |x|^{\frac13} +|v|$.
\end{theo}

\begin{rema}
The left hand side can be understood as a H\"older regularity of (kinetic) order $2+\alpha$, according the specific definition of H\"older spaces $\cC_\ell^\beta$, $\beta \ge 0$, given in the next section~(Definition~\ref{def:holder}). We compare this result to the classical Schauder estimate for parabolic equations in the next subsection.
\end{rema}

Such a Schauder estimate is typically used to reach well-posedness of nonlinear equations after the derivation of H\"older estimates on coefficients. In order to illustrate this fact, we consider in the second half of this paper the equation
\begin{equation}\label{eq:main}
\left(\partial_t + v \cdot \nabla_x\right) f = \rho[f] \nabla_v \cdot \left(\nabla_v f + v f \right)
\end{equation}
for an unknown $0 \le f =f(t,x,v)$, supplemented with the initial condition $f (0,x,v)\linebreak=f_{in}(x,v)$ in $\T^d \times \R^d$, where $\rho[f] (t,x):= \int_v f(t,x,v) \dd v$ and $\T^d$ denotes the $d$- dimensional torus. We emphasize the fact that studying~\eqref{eq:main} with $x \in \T^d$ is equivalent to study it with $x \in \R^d$ with periodic initial data. The known a priori estimates that are preserved in time for this equation are $L^1(\T^d \times \R^d)$ and $C_1 \mu \le f \le C_2 \mu$, where $\mu$ denotes the Gaussian $(2\pi)^{-d/2} e^{-|v|^2/2}$. They are not sufficient to derive uniqueness or bootstrap higher regularity. The Schauder estimate from Theorem~\ref{thm:variable}, together with the H\"older regularity from~\cite{gimv} (see Theorem~\ref{t:gimv}), allows us to prove global well-posedness of Eq.~\eqref{eq:main} in Sobolev spaces. In the following statement $H^k(\T^d \times \R^d)$ denotes the standard $L^2$-based Sobolev space.


\begin{theo}[Global well-posedness for a toy nonlinear
model]\label{theo:main-holder}
Given two constants $0<C_1\le C_2$, let $f_{in}$ be such that $f_{in}/\sqrt{\mu} \in H^k(\T^d \times \R^{d})$ with $k > 2+d/2$ and satisfying $ C_1 \mu \le f_{in} \le C_2 \mu$. There then exists a unique global-in-time solution $f$ of~\eqref{eq:main} in $(0,+\infty) \times \T^d \times \R^d$ satisfying $f(0,x,v) = f_{in} (x,v)$ everywhere in $\T^d \times \R^{d}$ and $f(t)/\sqrt{\mu} \in H^k(\T^d \times \R^{d})$ for all time $t>0$ and $C_1 \mu \le f \le C_2 \mu$ in $[0,+\infty) \times \T^d \times \R^d$.
\end{theo}



\subsection{Schauder estimates for kinetic equations}

The Schauder estimate for solutions $g(t,v)$ to parabolic equations of the form
\begin{multline}\label{eq:parabolic}
\partial_t g =
\sum_{1\,\le\:i,\,j\:\le\,d} a^{i,\,Sj} \partial^2_{v_iv_j} g +
\sum_{1\,\le\,i\,\le\,d}b^i \partial_{v_i} g + c g + S \quad \emph{ in } \R \times \R^d \\[1ex]
\text{takes the form } \; [\partial_t g]_{\cC^{\alpha}\left(\R\,\times\,\R^d\right)} + \left[D^2_v g\right]_{\cC^\alpha\left(\R\,\times\,\R^d\right)} \le C
\left([S]_{\cC^\alpha\left(\R\,\times\,\R^d\right)} + \|g\|_{L^\infty\left(\R\,\times\,\R^d\right)} \right)
\end{multline}
where $C^\alpha (\R \times \R^d)$ denotes the classical H\"older space with respect to the parabolic distance $\|(t,v)\|=|t|^{\frac12} + |v|$. This distance accounts for the parabolic scaling $(t,v)\linebreak\mapsto (r^2t,rv)$.

Because we work with kinetic Fokker--Planck equations, the usual parabolic scaling is replaced with the kinetic scaling $(t,x,v) \mapsto (r^2t,r^3x,rv)$. Moreover, parabolic equations of the form~\eqref{eq:parabolic} are translation invariant while kinetic Fokker--Planck equations of the form~\eqref{eq:linear} are translation invariant in the space variable $x$ but not in the velocity one $v$; the latter is replaced by the \emph{Galilean invariance}. As already noticed for instance in~\cite{gimv,polidoro1994}, these two facts ---kinetic scaling and Galilean invariance--- naturally require new definitions for cylinders, order of polynomials and H\"older continuity. We therefore define (see Definition~\ref{def:holder}) the space $\cC^\beta_\ell$ to be the set of functions whose difference with any polynomial of (kinetic) degree smaller than $\beta$ decays at rate $r^\beta$ in a cylinder of radius $r>0$. Following~\cite{imbert2018schauder}, one can define the \emph{kinetic degree} that follows the kinetic scaling, as $2 (\text{degree in } t) + 3(\text{degree in } x)\linebreak+ (\text{degree in } v)$. The subscript $\ell$ stands for ``left'' since the transformations leaving the equation invariant are applied to the left, see the next section.

There exists a well-developed literature of Schauder estimates for ultraparabolic equations, e.g.~\cite{manfredini,polidoro1994}, and so-called Kolmogorov or H\"ormander type equations, see e.g.~\cite{bb,dfp,lunardi1997,radkevich2008} and references there in. Some of these large classes of equations include the linear kinetic Fokker--Planck equations of the form~\eqref{eq:linear}. Moreover, \eqref{eq:linear} is already considered in~\cite{hs}. However in all these works, either the choice of H\"older spaces is not appropriate to the study of kinetic equations or the assumptions on the coefficients are too strong or the estimate is too weak.

In~\cite{hs}, the authors use the same natural H\"older spaces $\cC_\ell^\alpha$ for $\alpha \in (0,1)$ but make other choices for higher exponents $\alpha$. For instance, following the aforementioned classical Schauder estimate for parabolic equations, the semi-norm $[\cdot]_{2+\alpha,\,Q}$ in~\cite{hs} equals the sum of $[D^2u]_{\alpha,Q}$, $[\partial_t u]_{\alpha,\,Q}$ \emph{plus} an additional semi-norm controlling $x$-variations. Such a choice can be compared the equivalence of norms discussed in Remark~\ref{rem:equiv}. The authors of~\cite{hs} explain their choice by the fact that the natural Schauder estimate (in the spirit of Theorem~\ref{thm:variable}) only provides a regularity in the $x$ variable of order $(2+\alpha)/3 <1$ while they aim at reaching complete smoothing by bootstrap. Note that such bootstrap can nevertheless be achieved in our spaces $\cC_\ell^\beta$, but this requires to work with difference derivatives in $x$ of fractional order $(2+\alpha)/3$ each time the Schauder estimate is applied. Due to the technical length of this argument when providing full details and in order to keep this paper concise, we defer this higher regularity bootstrap on the nonlinear model~\eqref{eq:main} to a future study, and also note that such bootstrap techniques are being also implemented for the Boltzmann equation in~\cite{bootstrap}.

We also remark that our choice of norms ---that measures regularity by estimating the oscillations of a Taylor expansion remainder--- is related to the proof of the Schauder estimate. Indeed, we adapt the argument by Safonov~\cite{safonov1984} in the parabolic case, explained in Krylov's book~\cite{krylov1996}. In the latter argument, the oscillation of the remainder of the second-order Taylor expansion of the solution is shown to decay at rate $r^{2+\alpha}$ in a cylinder of radius $r$, and a corrector is introduced to the second-order Taylor polynomial to account for the contribution of the source term at large distance. Compared with the parabolic argument in~\cite{krylov1996}, the main conceptual difference is in the proof of the \emph{gradient bound}, see Proposition~\ref{prop:grad}. We combine Bernstein's method, as in~\cite{krylov1996}, with ideas and techniques borrowed from the hypocoercivity theory~\cite{villani2009memoir}.

\subsection{Motivation and background for the toy model}

Equation~\eqref{eq:main} describes the evolution of the probability density function $f$ of particles. The free transport translates the fact that the variable $v$ is the velocity of the particle at position $x$ at time $t$. The operator $\rho [f] \nabla_v\,\cdot\,(\nabla_v f + vf)$ takes into account the interaction between particles. The diffusion coefficient $\rho [f]$ is proportional to the total mass of particles lying at $x$ at time $t$: diffusion is strong in regions where local density is large and weak in regions where local density is small. The diffusion operator $\nabla_v\,\cdot\,(\nabla_v f + vf)$ is chosen so that the unique steady state of this equation is the Gaussian $\mu(v) = (2\pi)^{-d/2} e^{-|v|^2/2}$.

Even if it has a substantially simpler structure, the nonlinear kinetic equation~\eqref{eq:main} shares some similarities with the equation derived by Landau in 1936~\cite{landau1936} as a diffusive limit, in the regime when grazing collisions dominate, of the (Maxwell)--Boltzmann equation discovered in 1867--1872~\cite{boltzmann1872weitere,clerk1867dynamical} for describing rarefied gases. These latter equations are respectively referred to as the Landau (or Landau--Coulomb) equation and the Boltzmann equation. Landau derived this equation in order to describe plasmas made up of electrons and ions interacting by Coulombian forces. The Boltzmann collision operator can describe long-range interactions less singular than the Coulomb interactions but does not make sense for the Coulomb interaction. The Landau collision operator is of the form $\nabla_v\,\cdot\,(A_f \nabla_vf + \vec b_f f)$ with $A_f = A \star f$ and $\vec b_f = \vec b \star f$ ($\star$ stands for the convolution with respect to the velocity variable) where $A$ is the matrix $\frac{1}{|v|} (\Id -\frac{v \otimes v}{|v|^2})$ and $\vec b= - \nabla_v\,\cdot\,A$. The collision operator in equation~\eqref{eq:main} corresponds to the (much simpler) case where coefficients are given by $A_f = \rho[f] \Id$ and $B_f = \rho[f] v$.

In the case of the Landau equation, both coefficients are defined by integral quantities involving the solution. Our simplified toy model~\eqref{eq:main} replaces these convolutions crudely by their averages and neglects the issues of the various positive or negative moments at large velocities. This explains the factor $\rho[f]$. Our simplified toy model also shares the same Gaussian steady state as the Landau collision operator.

This simplification respects the principle at the source of nonlinearity in bilinearity collision operators: that the amount of collisions at a point is related to the local density of particles. Note that replacing $\rho[f]$ by another $v$-moment of the solution, or even having different $v$-moments in front respectively of the diffusion and drift terms, could most likely be treated by variants of the method developed in this paper. It is also likely that replacing $\rho[f]$ by $F(\rho[f])$ where $F:\R_+^* \to \R_+^*$ is a smooth nonlinear map could be treated by variants of the methods in this paper.

The model~\eqref{eq:main} was also studied in~\cite{MR2252155} (see~\cite[equation$\MK$(9)]{MR2252155}, when keeping only mass conservation) and the authors show how its spatially homogeneous version arise as a mean-field limit of an $N$-particle Markov process in the spirit of Kac's process~\cite{MR0084985}. It is also related to the gallery of nonlinear Fokker--Planck models discussed for instance in~\cite{chavanis2008nonlinear}.

A recent line of research consists in extending methods from the elliptic and parabolic theories to kinetic equations. Silvestre in particular made key progresses on the Boltzmann equation without cut-off in~\cite{MR3551261}, and together with the first author later obtained local H\"older estimate in~\cite{is} and a Schauder estimate for a class of kinetic equations with integral fractional diffusion in~\cite{imbert2018schauder}. In parallel, a similar program was initiated for the Landau equation in~\cite{gimv}, following up from an earlier result in~\cite{wz09}. The local H\"older estimate is obtained in~\cite{gimv,wz09} for essentially bounded solutions, the Harnack inequality is proved in~\cite{gimv} and some Schauder estimates are derived in~\cite{hs}. The results contained in the present article are part of this emerging trend in kinetic theory.

\subsection{Strategy of proof for global well-posedness}

We explain here the various ingredients used in the proof of the global well-posedness of equation~\eqref{eq:main}. The proof proceeds in 5 steps.
\begin{enumerate}
\item \label{1.4.1} First, the maximum principle implies that if the initial datum $f_{in}$ lies between $C_1 \mu$ and $C_2 \mu$, then the corresponding solution $f$ to~\eqref{eq:main} satisfies the same property: as long as the solution $f(t,x,v)$ is well-defined, we have $C_1 \mu (v) \linebreak\le f(t,x,v) \le C_2 \mu (v)$ for all $(t,x,v)$ (see Lemma~\ref{l:gaussian}). In particular, this ensures that the solution $f$ has fast decay at large velocities and that the diffusion coefficient $\rho [f]$ satisfies $C_1 \le \rho [f](t,x) \le C_2$ for all $(t,x)$. Therefore, the equation satisfies the uniform ellipticity condition in $v$ as stated in~\eqref{eq:elliptic}.
\item \label{1.4.2} We deduce from the bound on $\rho [f]$ that the solution $f$ satisfies an equation of the form
\[
(\partial_t + v \cdot \nabla_x) f = \nabla_v \cdot (A \nabla_v f+
\vec b f)
\]
for a symmetric real matrix $A$ whose eigenvalues all lie in $[C_1,C_2]$. In particular, we can use the local H\"older estimate from~\cite{gimv,wz09}, see Theorem~\ref{t:gimv} from Subsection~\ref{ss:holder}. The decay estimate from Step~1 and the H\"older regularity $\cC_\ell^{\alpha_0}$ for some small $\alpha_0$ are then combined with the Schauder estimate from Theorem~\ref{thm:variable} to derive a \emph{higher-order} H\"older estimate in $C^{2+\alpha_0}_\ell$ (see Proposition~\ref{prop:local}).
\item \label{1.4.3} With such a higher order H\"older estimate at hand, we next study how Sobolev norms in $x$ and $v$ grow as time increases and we derive a \emph{continuation criterion} (in the same spirit as the Beale--Kato--Majda blow-up criterion~\cite{MR763762}). We prove then that the blow-up is prevented by the $\cC_\ell^{2+\alpha_0}$ H\"older estimate from Step~\ref{step2}. This finally yields global well-posedness of equation~\eqref{eq:main} in Sobolev spaces.
\end{enumerate}

It is worth mentioning that a conditional global smoothing effect for the Landau equation with moderately soft potentials has been recently obtained in~\cite{hs} by combining the ingredients listed in Steps~\ref{step1} to Step~\ref{step3} above. Moreover, establishing such a global smoothing effect is in progress for the Boltzmann equation without cut-off with moderately soft potentials~\cite{bootstrap}. These works however assume a priori that some quantities such as mass, energy and entropy densities remain under control along the flow. The a priori assumption is necessary to prove, among other things, that the equations enjoy some uniform ellipticity and to establish good decay in the velocity variable. The interest of the toy nonlinear model~\eqref{eq:main} lies in the fact that it is a nontrivial and physically relevant model for which \emph{unconditional} global well-posedness can be proven following such a programme. The main simplification of our model compared with the Landau equation is the lack of local conservation of momentum and energy; therefore the fluid dynamics on the local density, momentum and energy fields reduces to the heat flow in the fluid limit and avoids the difficulties of the Euler and Navier--Stokes dynamics. The results of this paper hence provide one more hint that the formation of singularities, if any, in the Cauchy problem for non-linear kinetic equations is likely to come from~\eqref{1.4.1} fluid mechanics, or~\eqref{1.4.2} issues with the decay at large velocities.

\subsection{Perspectives}

We conclude the introduction by mentioning that the well-posedness result for the toy nonlinear model can be improved in two directions. First, more general initial data could be considered by constructing solutions directly in our H\"older spaces rather than mixing H\"older and Sobolev spaces. Second, we previously mentioned that $C^\infty$ regularization is expected for positive times by applying iteratively the Schauder estimates.

\subsection{Organisation of the article}

Section~\ref{sec:prelim} is devoted to the definition of H\"older spaces. The Schauder estimate from Theorem~\ref{thm:variable} is proved in Section~\ref{sec:schauder}. We prove Theorem~\ref{theo:main-holder} in the final Section~\ref{sec:toy} by constructing local solutions to the non-linear equation~\eqref{eq:main} in Sobolev spaces and by using the Schauder estimate to extend these solutions globally in time.

\subsection{Notation}
We collect here the main notations for the convenience of the reader.

\subsubsection{Euclidian space and torus}
The $d$-dimensional Euclidian space is denoted by $\R^d$ and the $d$-dimensional torus by $\T^d$. Throughout this article, the space variable $x$ belongs to $\R^d$, except in Section~\ref{sec:toy} where $x \in \T^d$.

\subsubsection{Multi-indices}
The \emph{order} of $m \in \N^d$ is $|m| := m_1+\dots+m_d$. Given a vector $x \in \R^d$ and $m \in \N^d$, we denote $x^m := \prod_{i=1} ^d x_i^{m_i}$.

\subsubsection{Balls and cylinders}
$B_r$ denotes the open ball of $\R^d$ of radius $r$ centered at the origin. $Q_r(z_0)$ denotes a cylinder in $\R \times \R^d \times \R^d$ centered at $z_0$ of radius $r$ following the kinetic scaling, see~\eqref{eq:def-cyl}. $Q_r$ simply denotes $Q_r((0,0,0))$. The scaled variable is $S_r (z):=(r^2t,r^3x,rv)$ for $z=(t,x,v)$, see~\eqref{eq:scaling}. Radii of cylinders are denoted by $r$. Unless further constraints are stated, this radius is an arbitrary positive real number. It is sometimes restricted by the fact that a cylinder should not leak out of the domain of study of the equation (in particular in time), sometimes chosen to be $1$ or $2$, or multiplied by a given constant, e.g. $3/2$ or $K+1$.

\subsubsection{Constants}
We use the notation $g_1 \lesssim g_2$ when there exists a constant $C > 0$ independent of the parameters of interest such that $g_1 \leq C g_2$. We analogously write $g_1 \gtrsim g_2$. We sometimes use the notation $g_1 \lesssim_{\delta} g_2$ if we want to emphasize that the implicit constant depends on some parameter $\delta$.

\subsubsection{H\"older spaces and exponents}
Given an open set $\cQ$, $\cC_\ell^\alpha (\cQ)$ denotes the set of $\alpha$-H\"older continuous functions in $\cQ$, see Definition~\ref{def:holder}. The subscript $\ell$ refers to ``left'', it is not a parameter. The letter $\alpha$ denotes an arbitrary positive exponent. This exponent will be fixed to some value $\alpha_0$ in Section~\ref{sec:toy} after applying the local H\"older estimate from~\cite{gimv,wz09} recalled in Subsubsection~\ref{ss:holder}.

\subsubsection{Kolmogorov operator}
The Green function of the operator $\cL_K := \partial_t + v \,\cdot\,\nabla_x-\Delta_v$ is denoted by $\cG$.

\section{H\"older spaces}\label{sec:prelim}

\subsection{Galilean invariance, scaling and cylinders}

Given $z := (t,x,v) \in \R \times \R^d \times \R^d$ and $r>0$ define
\begin{equation}\label{eq:scaling}
S_r(z) := \left(r^2 t,r^3 x, r v\right).
\end{equation}

The Galilean invariance of equations~\eqref{eq:linear} and~\eqref{eq:main} is expressed with the non-commutative product
\[
\left(t_1,x_1,v_1\right) \circ \left(t_2,x_2,v_2\right) := \left(t_1+t_2,x_1 +x_2 + t_2 v_1,v_1 + v_2\right)
\]
with inverse element denoted $z^{-1} := (-t,-x+tv,-v)$ for $z:=(t,x,v)$.

Given $z_0 \in \R \times \R^d \times \R^d$ and $r >0$, we define the unit cylinder $Q_1:=Q_1(0)\linebreak=(-1,0]\times B_1 \times B_1$, then $Q_r:=Q_r(0) = S_r(Q_1)$ and finally $Q_r (z_0) = \{ z_0 \circ z: z \in Q_r\}$. This results in
\begin{multline}\label{eq:def-cyl}
Q_r (z_0)\\
:= \left\{ (t,x,v) : t_0 - r^2 < t \le t_0, \, |x-x_0 - (t-t_0)v_0| < r^3, \, |v-v_0| < r \right\}.
\end{multline}

\subsection{The Green function}

Consider the Kolmogorov equation
\begin{equation}\label{eq:kolmogorov}
(\partial_t + v \cdot \nabla_x) g = \Delta_v g + S
\end{equation}
with a given source term $S$. The Green function $\cG$ of the operator $\cL_K:= \partial_t + v\,\cdot\,\nabla_x - \Delta_v$ was computed by Kolmogorov~\cite{kolm}:
\begin{equation}\label{eq:green}
\cG(z) =
\begin{cases}
\left(\frac{\sqrt{3}}{2 \pi t^2}\right)^d\,e^{-\frac{3\left|x+\frac{t}2 v\right|^2}{t^3}} e^{-\frac{|v|^2}{4 t}} & \text{ if } t > 0, \\
0 & \text{ if } t \le 0.
\end{cases}
\end{equation}

\begin{prop}[Properties of the Green function]\label{prop:pregreen}
Given $S \in L^\infty(\R \times \R^d \times \R^d)$ with compact support in time, the function
\begin{multline*}
g(t,x,v) \\
\begin{aligned}
&= \int_{\R\,\times\,\R^d\,\times\,\R^d}
\cG\left(\tilde z^{-1} \circ z\right) S(\tilde z) \dd \tilde t \dd
\tilde x \dd \tilde v \quad
 \left(\text{with } z:=(t,x,v) \text{ and } \tilde z := \left(\tilde t,\tilde x,\tilde v\right)\right)\\
& = \int_{\R\,\times\,\R^d\,\times\,\R^d} \cG\left(t-\tilde t,x- \tilde x - \left(t - \tilde t\right)\tilde v,v-\tilde v\right) S \left(\tilde t,\tilde x,\tilde v\right) \dd \tilde t \dd \tilde x \dd \tilde v =: (\cG \starkin S)(z)
\end{aligned}
\end{multline*}
satisfies~\eqref{eq:kolmogorov} in $\R \times \R^d \times \R^d$ (the ``kinetic'' convolution $\starkin$ follows the Galilean invariance).

Moreover, for all $z_0 =(t_0,x_0,v_0) \in \R \times \R^d \times \R^d$ and $r >0$
\[
\left\|\cG \starkin \un_{Q_r(z_0)}\right\|_{L^\infty\left(Q_r(z_0)\right)}
\lesssim_d r^2.
\]
\end{prop}

\begin{proof}
The argument of~\cite[Lemma~8.4.1, p.~115]{krylov1996} applies similarly: given $z \in Q_r(z_0)$, calculate
\begin{align*}
\cG \starkin \un_{Q_r(z_0)}(z) & = \int_{Q_r(z_0)} \cG\left(t-\tilde t,x- \tilde x - \left(t - \tilde t\right)\tilde v,v-\tilde v\right) \dd \tilde z \\
& = r^2 \int_{Q_1(z_0)} \cG\left(\frac{t}{r^2}-\bar t,\frac{x}{r^3}- \bar x - \left(\frac{t}{r^2} - \bar t\right)\bar v,\frac{v}{r}-\bar v\right) \dd \bar z \\
& = r^2 \cG \starkin \un_{Q_1(z_0)}\left(S_\frac1r(z)\right)
\end{align*}
which yields the result.
\end{proof}

\subsection{H\"older spaces}

We now introduce H\"older spaces similar to that in~\cite{imbert2018schauder}.
\begin{defi}[H\"older spaces]\label{def:holder}
The \emph{kinetic degree} of a monomial $m (t,x,v) = t^{k_0} \Pi_{i=1}^d x_i^{k_i} v_i^{l_i}$ associated with the \emph{partial} degrees $k_0,k_1,\,\dots,\,k_d \in \N$ and $l_1,\,\dots,\,l_d \in \N$ is defined as
\[
\degkin m := 2 k_0 + 3 \left(\sum_{i=1}^d k_i\right) + \sum_{i=1}^d l_i.
\]
In particular, constants have zero kinetic degree. The kinetic degree $\degkin p$ of a polynomial $p \in \R[t,x,v]$ is defined as the largest kinetic degree of the monomials $m_j$ appearing in $p$.

Given an open set $\cQ \subset \R \times \R^d \times \R^d$ and $\beta >0$, we say that a function $g: \cQ \to \R$ is $\beta$-H\"older continuous at a point $z_0 \in \cQ$ if there is a polynomial $p \in \R [t,x,v]$ with $\degkin p < \beta$ and a constant $C>0$ such that
\begin{equation}\label{e:p}
\forall \ r >0, \quad \|g- p\|_{L^\infty\left(Q_r\,(z_0)\,\cap\,\cQ\right)} \le C r^{\beta}.
\end{equation}
If this property holds true for all $z_0 \in \cQ$, the function $g$ is $\beta$-H\"older continuous in $\cQ$ and we write $g \in \cC_\ell^\beta (\cQ)$. The smallest constant $C$ such that the property~\eqref{e:p} holds true for all $z_0 \in \cQ$ is denoted by $[g]_{\cC^{\beta}(\cQ)}$.

The $\cC_\ell^{\beta}$-norm of $g$ is then $\|g\|_{\cC_\ell^{\beta}(\cQ)} := \|g\|_{L^\infty(\cQ)} + [g]_{\cC_\ell^{\beta}(\cQ)}$.
\end{defi}
\goodbreak
\begin{rema}\ \\*[-1em]
\begin{enumerate}
\item \label{rema2.3.1} For $\beta \in (0,1)$, the semi-norm $\cC^\beta_\ell$ is equivalent to the standard H\"older semi-norm $\cC^\beta$ for the distance $\|(t,x,v)\|=|t|^{\frac12} + |x|^{\frac13} + |v|$.

\item \label{rema2.3.2}For a non-zero integer $k \in \N$, the spaces $C^k_\ell$ differ from the usual $C^k$ spaces, in that the highest-order derivatives are not continuous but merely $L^\infty$. For instance $C^1_\ell$ functions are Lipschitz continuous in $v$ but not continuously differentiable in $v$.


\item \label{rema2.3.3}In~\cite{imbert2018schauder}, a Schauder estimate is obtained for kinetic equations with an integral collision operator with fractional ellipticity in the velocity variable $v$. A parameter $s \in (0,1)$, related to the order of differentiation of the integral operator, plays a role in the definition of cylinders (through scaling) and, in turn, in the definition of kinetic degree and H\"older spaces. The present definition corresponds to the case $s=1$.

\item \label{rema2.3.4}The subscript $\ell$ emphasizes the fact that this definition is based on the non-commutative product $\circ$ defined above. The symbol $\ell$ stands for ``left'': because of Galilean invariance, the kinetic equations considered in~\cite{imbert2018schauder} and the present work are left invariant with respect to this product. See~\cite{imbert2018schauder} for further discussion.

\item \label{rema2.3.5}In~\cite{imbert2018schauder}, H\"older spaces are defined by using a kinetic distance $d_\ell$. It is pointed out that this distance satisfies $\frac14 \|z_2^{-1} \circ z_1\| \le d_\ell (z_1,z_2) \le
\|z_2^{-1} \circ z_1\|$ where $\|(t,x,v)\|=|t|^{\frac12} + |x|^{\frac13} + |v|$. This inequality justifies the fact that the Definition~\ref{def:holder} above coincides with~\cite[Definition~2.3]{imbert2018schauder} and that semi-norms only differ by a factor $4^\beta$.

\item \label{rema2.3.6}When $\cQ = \R \times \R^d \times \R^d$, we simply write $[\cdot]_{\cC_\ell^{\beta}}$, $\|\cdot\|_{\cC_\ell^{\beta}}$, $\|\cdot\|_{L^\infty}$, etc.
\end{enumerate}
\end{rema}

\subsection{Second order Taylor expansion}

When $\beta =2+\alpha \in (2,3)$, we now prove that the polynomial $p$ realizing the infimum in the $\cC_\ell^\beta$-semi-norm is the Taylor expansion of kinetic degree~$2$:
\begin{multline}\label{e:taylor}
\cT_{z_0}[g](t,x,v)
 := g(z_0) + (t-t_0) \left[\partial_t g + v_0
\cdot \nabla_x g \right](z_0)\\
+ (v-v_0) \cdot \nabla_v g(z_0) + \frac12 (v-v_0)^T \cdot D^2 _v g(z_0) \cdot (v-v_0).
\end{multline}
Remark that the linear part in $x$ does not appear since it is of kinetic degree~$3$.

We recall and denote
\begin{align*}
[g]_{\cC_\ell^{2+\alpha}(\cQ)} &\; = \sup \left\{ \frac{ \inf_{p\,\in\,\bbP}
\|g-p \|_{L^\infty\left(Q_r(z_0) \cap \cQ\right)}}{r^{2+\alpha}} : z_0 \in \cQ, \ r >0 \right\} \\
[g]_{\cC_{\ell,\,0}^{2+\alpha}(\cQ)} & := \sup \left\{\frac{\|g-\cT_{z_0}[g]
\|_{L^\infty\left(Q_r(z_0) \cap \cQ\right)}}{r^{2+\alpha}}\quad : z_0 \in \cQ, \ r >0 \right\}
\end{align*}
where $\bbP$ denotes the set of polynomials of kinetic degree smaller than or equal to $2$.


\begin{lemm}\label{lem:NN0}
Given $\alpha \in (0,1)$, there exists $C\in (0,1)$ such that for all $g \in \cC_\ell^{2+\alpha}(\cQ)$, the derivatives in $\cT_{z_0}[g]$ exist for all $z_0 \in \cQ$ and
\[
C [g]_{\cC_{\ell,\,0}^{2+\alpha}(\cQ)} \le [g]_{\cC_\ell^{2+\alpha}(\cQ)} \le [g]_{\cC_{\ell,\,0}^{2+\alpha}(\cQ)}.
\]
\end{lemm}


\begin{proof}
First reduce to $z_0=0$ by the change of variables $g^\sharp(z) := g(z_0 \circ z)$. We continue however to simply call the function $g$. We start by proving the first inequality. One needs to identify the minimizer $p \in \R[t,x,v]$ realising $\inf_{p\,\in\,\R[t,\,x,\,v]} \|g-p \|_{L^\infty(Q_r)}$ in the limit $r \to 0^+$. Let $\eps >0$ and consider $r_k = 2^{-k}$ and $p_k \in \bbP$ s.t. $\|g-p_k \|_{L^\infty(Q_{r_k})}\linebreak\le r_k^{2+\alpha}([g]_{\cC_\ell^{2+\alpha}(\cQ)} + \eps)$. Write $p_k(t,v) =: a_k + b_k t + q_k \cdot v + \frac12 v^T M_k v$. By subtraction one gets
\begin{align*}
\left\| p_{k+1}-p_{k} \right\|_{L^\infty\left(Q_{r_{k+1}}\right)} \le 2 r_k^{2+\alpha} \left([g]_{\cC_\ell^{2+\alpha}(\cQ)}+\eps \right)
\end{align*}
which writes in terms of the coefficients
\begin{align*}
\left\| \left(a_{k}-a_{k+1}\right) + \left(b_{k} - b_{k+1}
\right) t + \left(q_{k} - q_{k+1} \right) \cdot v + \frac12 v^T
\left(M_{k} - M_{k+1} \right) v
\right\|_{L^\infty\left(Q_{r_{k+1}}\right)} \\
\le 2 r_{k}^{2+\alpha} \left([g]_{\cC_\ell^{2+\alpha}(\cQ)} + \eps \right).
\end{align*}


Testing for $t=0$ and $v=0$ gives $|a_{k}-a_{k+1}| \lesssim r_{k}^{2+\alpha} ([g]_{\cC_\ell^{2+\alpha}(\cQ)}+\eps)$. Using the latter and testing for $v=0$ and $|t|=r_{k+1}^2$ gives
\[
\left|b_{k} - b_{k+1}\right| \lesssim r_{k}^{\alpha} \left([g]_{\cC_\ell^{2+\alpha}(\cQ)}+\eps\right).
\]
Testing for $t=0$ and summing $v$ and $-v$ with $|v|=r_{k+1}$ in all directions gives $|M_{k} - M_{k+1}| \lesssim r_{k}^{\alpha} ([g]_{\cC_\ell^{2+\alpha}(\cQ)}+\eps)$. Finally by difference and testing with $t=0$ and all directions of $|v|=r_k$, one gets
\[
\left|q_{k} - q_{k+1}\right| \lesssim r_{k}^{1+\alpha} \left([g]_{\cC_\ell^{2+\alpha}(\cQ)}+\eps\right)\,.
\]
This shows that the coefficients are converging with
\begin{align*}
\left|a_{k}-a_{\infty}\right| &\lesssim r_{k}^{2+\alpha} \left([g]_{\cC_\ell^{2+\alpha}(\cQ)}+\eps\right)\text{ and }\left|b_{k} - b_{\infty}\right| \lesssim r_{k}^{\alpha} \left([g]_{\cC_\ell^{2+\alpha}(\cQ)}+\eps\right)\\ 
\text{ and }\left|q_{k} - q_{\infty}\right| &\lesssim r_{k}^{1+\alpha} \left([g]_{\cC_\ell^{2+\alpha}(\cQ)}+\eps\right)
\text{ and }\left|M_{k} - M_{\infty}\right| \lesssim r_{k}^{\alpha} \left([g]_{\cC_\ell^{2+\alpha}(\cQ)}+\eps\right).
\end{align*}
These convergences and estimates imply that
\[
\left\| g - P_\infty \right\|_{L^\infty\left(\cQ_{r_k}\right)} \lesssim r_{k}^{2+\alpha} \left([g]_{\cC_\ell^{2+\alpha}(\cQ)}+\eps\right)
\]
which in turn implies that $a_\infty = g(0,0,0)$ and $b_\infty = \partial_t g(0,0,0)$ and $q_\infty = \nabla_v g(0,0,0)$ and $M_\infty = D^2 _v g(0,0,0)$ (and proves the existence of such derivatives). We thus proved that
\[
\left\|g-\cT_0[g] \right\|_{L^\infty\left(Q_{r_k}\right)} \lesssim r_{k}^{2+\alpha} \left([g]_{\cC_\ell^{2+\alpha}(\cQ)}+\eps\right)
\]
where the constant does not depend on $k$. This in turn implies that same inequality for any $r>0$, with a constant at most multiplied by $2$, which concludes the proof since $\eps$ is arbitrarily small.

The proof of the second inequality $[g]_{\cC_\ell^{2+\alpha}(\cQ)} \le [g]_{\cC_{\ell,\,0}^{2+\alpha}(\cQ)}$ then follows from the existence of the derivatives appearing in $\cT_0[g]$ showed in the previous step, and the fact that $\cT_{0}[g]$ is of kinetic degree strictly smaller than $2+\alpha$.
\end{proof}



\subsection{Interpolation inequalities}

Interpolation inequalities are needed in the proof of the Schauder estimate.

\begin{lemm}[Interpolation inequalities]\label{l:interpol}
Let $\Qint, \Qext$ be two cylinders of the form $Q_\rho(z_c)$ and $Q_R(z_c)$ with either $\rho < R$ or $\rho = R = +\infty$. There exist $C,\beta>0$ depending only on $d$, $\alpha \in (0,1)$ (and $R-\rho$ if $R$ is finite) such that for any $\eps \in (0,1)$ and any $g \in \cC_\ell^{2+\alpha}(\Qext)$,
\begin{align}\label{eq:interp1}
[g]_{\cC^\alpha_\ell (\Qint)} & \le \eps [g]_{\cC_\ell^{2+\alpha}(\Qext)} + C \eps^{-\beta} \|g\|_{L^\infty(\Qext)} \\\label{eq:interp3}
\left[\nabla_v g\right]_{\cC_\ell^{\alpha}(\Qint)} & \le \eps [g]_{\cC_\ell^{2+\alpha} (\Qext)} + C\eps^{-\beta} \|g\|_{L^\infty (\Qext)}\\\label{eq:interp4-bis}
\left[D^2 _v g\right]_{\cC_\ell^\alpha\left(\R^{2d+1}\right)} & \le C [g]_{\cC_\ell^{2+\alpha}\left(\R^{2d+1}\right)} \\\label{eq:interp5-bis}
\left[\left(\partial_t + v \cdot \nabla_v\right) g\right]_{\cC_\ell^\alpha\left(\R^{2d+1}\right)} & \le C [g]_{\cC_\ell^{2+\alpha}\left(\R^{2d+1}\right)}.
\end{align}
Note that the two last inequalities are only proved in the case of the whole space $\rho=R=\infty$.
\end{lemm}


\begin{rema}
A similar result was also proved in~\cite[Lemma~2.7]{imbert2018schauder} by a different argument and in the terminology of \emph{kinetic degree}, see Definition~\ref{def:holder}. Denote $D=\Id,\nabla_v, D^2_v,(\partial_t + v \cdot \nabla_x)$ and define $k \in \{0,1,2\}$ the kinetic order of the differential operator $D$, i.e. $k=0$ for $D=\Id$, $k=1$ for $D = \nabla_v$ and $k=2$ for $D=D^2_v$ and $D=\partial_t + v \cdot \nabla_x$ (this is the kinetic degree of the polynomial naturally associated with the differential operator). Then given for $\beta \in \{0,\alpha\}$, and provided that $k+\beta < 2+\alpha$, we have
\begin{equation}\label{e:interpol-cyril}
\|Dg\|_{\cC_\ell^\beta(\cQ)} \le \eps^{(2+\alpha)-(k+\beta)} [g]_{\cC_\ell^{2+\alpha}(\cQ)} + C\eps^{-(k+\beta)} \|g\|_{L^\infty(\cQ)}
\end{equation}
where constant $C$ only depends on dimension $d$, $\alpha$ and $\|\cdot\|_{L^\infty(\cQ)}$. Observe that the restrictions imposed on $k,\beta$ yield the same inequalities as in Lemma~\ref{l:interpol}. The two key steps in~\cite[Lemma~2.7]{imbert2018schauder} are the proof of the general interpolation inequality
\[
[g]_{\cC_\ell^{k+\beta}(\cQ)} \le
\|g\|_{L^\infty(\cQ)}^{1-\frac{k+\beta}{2+\alpha}} [g]_{\cC_\ell^{2+\alpha}(\cQ)}^{\frac{k+\beta}{2+\alpha}} + C
\|g\|_{L^\infty(\cQ)}
\]
and inequalities relating the H\"older semi-norms of derivatives of a given function to its (higher order) semi-norm, as in~\eqref{eq:interp4-bis}-\eqref{eq:interp5-bis} but more general.
\end{rema}


\begin{rema}
It is possible to get rid of the condition $\rho < R$ when $R<+\infty$ but this requires substantial modifications. In the following proof, a global estimate ($\rho=R=+\infty$) is derived and a local ``interior'' one is easily obtained by a localization procedure. Since reaching $\rho =R<+\infty$ is irrelevant for this work and the proof below is different from the one contained in~\cite{imbert2018schauder}, we believe it can be useful to restrict ourselves to this special case.
\end{rema}

\begin{proof}[Proof of Lemma~\ref{l:interpol}]
It is sufficient to prove the interpolation inequalities~\eqref{eq:interp1}-\eqref{eq:interp3} in the case where $\Qint = \Qext = \R^{2d+1}$. Indeed, it is then sufficient to apply global estimates to $\tilde g = g \phi$ for some cut-off function $\phi$ such that $0 \le \phi \le 1$, $\phi =1$ in $\Qint$ and $\phi=0$ outside $\Qext$ to get the local ones. Indeed, in the case of~\eqref{eq:interp1},
\begin{align*}
[g]_{\cC^\alpha_\ell (\Qint)} & \le [g \phi]_{\cC^\alpha_\ell \left(\R^{2d+1}\right)} \\
& \le \eps [g \phi]_{\cC_\ell^{2+\alpha}\left(\R^{2d+1}\right)} + C \eps^{-\beta} \|g \phi\|_{L^\infty\left(\R^{2d+1}\right)} \\
& \le \eps \left([g]_{\cC_\ell^{2+\alpha}(\Qext)} + \|g\|_{L^\infty (\Qext)} [\phi]_{\cC_\ell^{2+\alpha}\left(\R^{2d+1}\right)}\right) + C \eps^{-\beta} \|g\|_{L^\infty (\Qext)} \\
& \le \eps [g]_{\cC_\ell^{2+\alpha}(\Qext)} + \tilde{C} \eps^{-\beta} \|g\|_{L^\infty(\Qext)}
\end{align*}
where in the case $0<\rho<R<+\infty$, the constant $\tilde{C} \sim (R-\rho)^{-2-\alpha}$ depends on the difference of radii. This localization is however not used in the proof of~\eqref{eq:interp4-bis}-\eqref{eq:interp5-bis} because it would add a term $\|g\|_{L^\infty(\Qext)}$ on the right hand side, which we want to avoid for the application of Remark~\ref{rem:equiv}.

Since we now work in the whole space $\R^{2d+1}$, we do not specify the domain of functions spaces in the remainder of the proof.

We next notice that a scaling argument allows us to only prove the inequalities with $\eps =1$: scaling then shows that $\beta = - \frac{\alpha}{2}$ in~\eqref{eq:interp1} and $\beta=-1-\alpha$ in~\eqref{eq:interp3}.
%\medskip
%{\bf  $L^\infty$ bounds on derivatives.} 
\begin{step}{$L^\infty$ bounds on derivatives\label{step1}.}

One can perform the same argument as in the proof of the preceding Lemma~\ref{lem:NN0} with $\eps = \| g \|_{L^\infty}$, $a_0=0$, $b_0=0$, $M_0=0$ as admissible first terms in the sequence to get
\begin{align*}
& \left|a_{k}-a_{\infty}\right| \lesssim r_{k}^{2+\alpha}
\left([g]_{\cC_\ell^{2+\alpha}}+ \| g \|_{L^\infty}\right) \\
& \left|b_{k} - b_{\infty}\right| \lesssim r_{k}^{\alpha}
\left([g]_{\cC_\ell^{2+\alpha}}+ \| g \|_{L^\infty}\right) \\
& \left|q_{k} - q_{\infty}\right| \lesssim r_{k}^{1+\alpha}
\left([g]_{\cC_\ell^{2+\alpha}}+ \| g \|_{L^\infty} \right) \\
& \left|M_{k} - M_{\infty}\right| \lesssim r_{k}^{\alpha}
\left([g]_{\cC_\ell^{2+\alpha}}+ \| g \|_{L^\infty} \right).
\end{align*}
This yields
\[
\left|g(z_0)\right|, \,\left|\partial_t g(z_0)+v\,\cdot\, \nabla_x g(z_0)\right|, \, \left|\nabla_v g(z_0)\right|, \, \left|\nabla_v^2 g(z_0)\right|
\lesssim [g]_{\cC_\ell^{2+\alpha}}+ \| g \|_{L^\infty}
\]
for any $z_0 \in \Qext$, which yields
\begin{align}\label{eq:interp2}
\left\|\nabla_v g\right\|_{L^\infty} & \lesssim [g]_{\cC_\ell^{2+\alpha}} + \|g\|_{L^\infty} \\\label{eq:interp4}
\left\| D^2 _v g \right\|_{L^\infty} & \lesssim [g]_{\cC_\ell^{2+\alpha}} + \|g\|_{L^\infty} \\\label{eq:interp5}
\left\| \left(\partial_t + v \cdot \nabla_v\right) g \right\|_{L^\infty} & \lesssim [g]_{\cC_\ell^{2+\alpha}} + \|g\|_{L^\infty}.
\end{align}
\end{step}

%\noindent {\bf Step~2. Control of lower-order H\"older norms.}
\begin{step}{Control of lower-order H\"older norms\label{step2}.}

 Using the previous $L^\infty$ bounds and Lemma~\ref{lem:NN0},
\begin{align*}
\left\|g-g(z_0)\right\|_{L^\infty(Q_r(z_0))} & \lesssim
\left\|g-\cT_{z_0}[g]\right\|_{L^\infty\left(Q_r(z_0)\right)} \\
+ & r \left(\left\| \partial_t g(z_0)+v \cdot \nabla_x g(z_0)
\right\|_{L^\infty} +
\left\| \nabla_v g(z_0) \right\|_{L^\infty} +
\left\|\nabla_v^2 g(z_0)\right\|_{L^\infty} \right) \\
& \lesssim r^{2+\alpha} [g]_{\cC^{2+\alpha}_\ell} + r \left([g]_{\cC_\ell^{2+\alpha}}+ \| g \|_{L^\infty} \right)
\end{align*}
which yields~\eqref{eq:interp1} with $\eps=1$.
\end{step}

%\noindent {\bf Step~3. H\"older regularity of first-order derivatives in $v$.}
\begin{step}{H\"older regularity of first-order derivatives in $v$.\label{step3}}

We are left with proving~\eqref{eq:interp3}, that is to say that for all $z_0 \in \R^{2d+1}$ and $z_1 \in Q_r (z_0)$, we have
\[
\left|\nabla_v g (z_1) - \nabla_v g (z_0) \right| \lesssim r^\alpha \left([g]_{\cC_\ell^{2+\alpha} } + \|g\|_{L^\infty }\right).
\]
In view of~\eqref{eq:interp4}, it is enough to consider $r \in (0,1]$.

Define for $z \in \R^{2d+1}$ and $u \in \bbS^{d-1}$ and $r \in (0,1]$,
\[
\sigma_{r,\,u}^1[g](z) := \frac{1}{r} \left[g(z \circ (0,0,ru)) -g(z) \right].
\]
We make two observations about this function. First there exists $\theta = \theta (z,r,u) \in (0,1)$ such that
\[
\sigma_{r,\,u} ^1[g] (z) = \nabla_v g (z \circ (0,0,\theta r u)) \cdot u.
\]
Second, for $p = \cT_{z_0} g(z)$, we have
\[
\sigma_{r,\,u}^1 [p](z_1) - \sigma_{r,\,u}^1 [p](z_0) = r u^T D^2_v g (z_0) (v_1-v_0).
\]

Let $z_0 \in \R^{2d+1}$, $z_1 \in Q_r (z_0)$ and $u \in \bbS^{d-1}$ be fixed, and $\delta >0$ to be chosen later. There then exists $\tilde z_0 \in Q_{\delta r} (z_0)$ and $\tilde z_1 \in Q_{\delta r} (z_1)$ such that
\[
\sigma_{\delta r,\,u}^1 [g] (z_i) = \nabla_v g (\tilde z_i) \cdot u, \qquad \text{ for } i=0,1.
\]
Then we can write, using Lemma~\ref{lem:NN0} to get the last inequality,
\begin{multline*}
\left| \left(\nabla_v g(z_1) - \nabla_v g(z_0)\right) \cdot u\right|
\le \left| \nabla_v g(z_1)\cdot u - \sigma _{\delta r,\, u}^1 [g](z_1) \right|\\
\begin{aligned}
 &\quad+ \left| \sigma _{\delta r,\,u}[g](z_1) - \sigma _{\delta r,\,u}^1 [p](z_1) \right| + \left| \sigma_{\delta r,\,u}^1 [p](z_1) - \sigma_{\delta r,\,u}^1 [p](z_0) \right| \\
&\quad +\left| \sigma_{\delta r,\,u}^1 [p](z_0) - \sigma_{\delta r,\,u}^1 [g](z_0) \right| + \left| \nabla_v g(z_0) \cdot u -\sigma_{\delta r,\,u}^1 [g](z_0) \right| \\
\le & \left| \nabla_v g(z_1) - \nabla_v g (\tilde z_1) \right| +\left\| \nabla^2 _v g \right\|_{L^\infty} |v_1-v_0| +\left| \nabla_v g(z_0) -\nabla_v g (\tilde z_0) \right| \\
&\quad + r^{-1} |(g-p) (z_0 \circ (0,0,\delta r u))| + r^{-1} |(g-p) (z_0)| \\
&\quad + r^{-1} |(g-p) (z_1 \circ (0,0,\delta r u))| + r^{-1} |(g-p) (z_1)| \\
\le & 2 (\delta r)^\alpha [\nabla_v g]_{\cC_\ell^\alpha} + r \left\| \nabla^2 _v g \right\|_{L^\infty} + 4 r^{1+\alpha} [g]_{\cC^{2+\alpha}_\ell }.
\end{aligned}
\end{multline*}
Using the fact that $r \le 1$ and taking the supremum over $z_0,z_1, u$, we get,
\[
\left[\nabla_v g\right]_{\cC_\ell^\alpha} \le 2 \delta ^\alpha \left[\nabla_v g\right]_{\cC_\ell^\alpha} + C \left(\left\| \nabla^2 _v g \right\|_{L^\infty} + [g]_{\cC^{2+\alpha}_\ell } \right).
\]
We can now pick $\delta >0$ such that $2 \delta^\alpha = \frac12$ and conclude that~\eqref{eq:interp3} holds true.
\end{step}


%\noindent {\bf Step~4. H\"older regularity of second-order derivative in $v$.}
\begin{step}{H\"older regularity of second-order derivative in $v$.\label{step4}}

We now prove~\eqref{eq:interp4-bis} following a similar strategy as in the previous step: we prove for all $z_0 \in \R^{2d+1}$ and $z_1 \in Q_r (z_0)$ with $r \in (0,1]$
\[
\left| \nabla_v ^2 g (z_1) - \nabla_v^2 g (z_0) \right| \lesssim r^\alpha [g]_{\cC_\ell^{2+\alpha}}.
\]

Define for $z \in \R^{2d+1}$ and $u \in \bbS^{d-1}$ and $r \in (0,1]$,
\[
\sigma^2 _{r,\,u}[g](z) := \frac{1}{r^2} \left[g(z \circ (0,0,ru)) + g(z \circ (0,0,-ru)) - 2 g(z) \right].
\]
This function satisfies for some $\theta = \theta (z,r,u) \in (-1,1)$ (and recalling $p=\cT_{z_0}[g]$)
\[
\sigma^2_{r,\,u} [g] (z) = u^T \nabla_v ^2 g (z \circ (0,0,\theta r u)) u \quad \text{ and } \quad \sigma_{r,\,u}^2[p](z_1) =
\sigma_{r,\,u}^2[p](z_0).
\]

Let $z_0 \in \R^{2d+1}$, $z_1 \in Q_r (z_0)$ and $u \in \bbS^{d-1}$ be fixed, and $\delta >0$ to be chosen later. There then exists $\tilde z_0 \in Q_{\delta r} (z_0)$ and $\tilde z_1 \in Q_{\delta r} (z_1)$ such that
\[
\sigma_{\delta r,\,u}^2 [g] (z_i) = u^T \nabla_v ^2g (\tilde z_i)u, \quad \text{ for } i=0,1.
\]
Then we can write, using Lemma~\ref{lem:NN0} to get the last inequality,
\begin{multline*}
\left| u^T \left(\nabla_v ^2 g(z_1) - \nabla_v ^2 g(z_0) \right)u\right|
\le \left| u^T \nabla_v ^2 g(z_1) u - \sigma _{\delta r,\,u}^2[g](z_1) \right| \\
\begin{aligned}
&\quad+ \left| \sigma _{\delta r,\,u}^2 [g](z_1) - \sigma _{\delta r,u}^2 [p](z_1) \right| + \left| \sigma_{\delta r,\,u}^2 [p](z_1) - \sigma_{\delta r,\,u}^2[p](z_0) \right| \\
&\quad+\left| \sigma_{\delta r,u}^2[p](z_0) - \sigma_{\delta r,\,u}^2 [g](z_0) \right| + \left| u^T \nabla_v ^2 g(z_0) u -\sigma_{\delta r,\,u}^2 [g](z_0) \right| \\
& \le \left| \nabla_v ^2 g(z_1) - \nabla_v ^2 g (\tilde z_1) \right| + + \left| \nabla_v ^2 g(z_0) -\nabla_v ^2 g (\tilde z_0) \right| \\
& \quad + r^{-2} |(g-p) (z_0 \circ (0,0,\delta r u))| + r^{-2} |(g-p) (z_0 \circ (0,0,-\delta r u))| \\
&\quad + 2 r^{-2} |(g-p) (z_0)| + r^{-2} |(g-p) (z_1 \circ (0,0,\delta r u))| \\
 & \quad + r^{-2} |(g-p) (z_1 \circ (0,0,-\delta r u))| + 2 r^{-2} |(g-p) (z_1)|\\
& \le 2 (\delta r)^\alpha \left[\nabla_v ^2g\right]_{\cC_\ell^\alpha} + 6 r^{\alpha} [g]_{\cC^{2+\alpha}_\ell }.
\end{aligned}
\end{multline*}
Using the fact that $r \le 1$ and taking the supremum over $z_0,z_1, u$, we get,
\[
\left[\nabla_v ^2 g\right]_{\cC_\ell^\alpha} \le 2 \delta ^\alpha \left[\nabla_v ^2 g\right]_{\cC_\ell^\alpha} + C [g]_{\cC^{2+\alpha}_\ell }
\]
which proves the inequality by choosing $2 \delta^\alpha = \frac{1}{2}$.
\end{step}
%
%\medskip

%\noindent {\bf Step~5. H\"older regularity of first-order transport derivative.}
\begin{step}{H\"older regularity of first-order transport derivative.\label{step5}}

The proof of~\eqref{eq:interp5-bis} follows the same strategy: we prove, denoting $Y:=\partial_t + v \cdot \nabla_x$, for all $z_0 \in \R^{2d+1}$ and $z_1 \in Q_r (z_0)$ with $r \in (0,1]$
\[
\left| Y g (z_1) - Y g (z_0) \right| \lesssim r^\alpha [g]_{\cC_\ell^{2+\alpha}}.
\]

Define for $z \in \R^{2d+1}$ and $r \in (0,1]$,
\[
\sigma^3 _{r}[g](z) := \frac{1}{r^2} \left[g(z \circ (r^2,0,0)) - g(z) \right].
\]
This function satisfies for some $\theta = \theta (z,r) \in (0,1)$
\[
\sigma^3_{r} [g] (z) = Y g (z \circ (\theta r^2,0,0) \quad \text{ and } \quad \sigma_{r}^3[p](z_1) =
\sigma_{r}^3[p](z_0).
\]

Let $z_0 \in \R^{2d+1}$, $z_1 \in Q_r (z_0)$, and $\delta >0$ to be chosen later. Then for some $\tilde z_0 \in Q_{\delta r} (z_0)$, $\tilde z_1 \in Q_{\delta r} (z_1)$:
\[
\sigma_{\delta r,u}^3 [g] (z_i) = Y g (\tilde z_i), \;\text{ for }\; i=0,1.
\]
Then we can write, using Lemma~\ref{lem:NN0} to get the last inequality,
\begin{align*}
\left| Y g(z_1) - Y g(z_0) \right| & \le \left| Y g(z_1) - \sigma _{\delta r}^3[g](z_1) \right|\\
&\quad+ \left| \sigma _{\delta r}^3 [g](z_1) - \sigma _{\delta r}^3 [p](z_1) \right| + \left| \sigma_{\delta r}^3 [p](z_1) - \sigma_{\delta r}^3[p](z_0) \right| \\
& \quad +\left| \sigma_{\delta r}^3[p](z_0) - \sigma_{\delta r}^3 [g](z_0) \right| + \left| Y g(z_0) -\sigma_{\delta r}^3 [g](z_0) \right| \\
& \le \left| Y g(z_1) - Y g (\tilde z_1) \right| +
\left| Y g(z_0) - Yg (\tilde z_0) \right| \\
& \quad + r^{-2} |(g-p) (z_0 \circ (\delta^2 r^2,0,0))| + r^{-2} |(g-p) (z_0)| \\
& \quad + r^{-2} |(g-p) (z_1 \circ (\delta^2 r^2,0,0))| + r^{-2} |(g-p) (z_1)| \\
& \le 2 (\delta r)^\alpha [Y g]_{\cC_\ell^\alpha} + 4 r^{\alpha} [g]_{\cC^{2+\alpha}_\ell}.
\end{align*}
Using the fact that $r \le 1$ and taking the supremum over $z_0,z_1, u$, we get,
\[
[Y g]_{\cC_\ell^\alpha} \le 2 \delta ^\alpha [Yg]_{\cC_\ell^\alpha} + C [g]_{\cC^{2+\alpha}_\ell }
\]
which proves the inequality by choosing $2 \delta^\alpha = \frac{1}{2}$.
\end{step}

This achieves the proof of the Lemma~\ref{l:interpol}.
\end{proof}


We will see below that~\eqref{eq:interp4-bis} and~\eqref{eq:interp5-bis} can be combined with the hypoelliptic estimate contained in Lemma~\ref{l:hypo} below to derive an equivalent semi-norm for the kinetic H\"older space $\cC_\ell^{2+\alpha}$, see Remark~\ref{rem:equiv}.

\subsection{A hypoelliptic estimate}

We investigate here how to recover the fact that a given function $g$ lies in $\cC^{2+\alpha}_\ell$ only knowing that its free transport and its velocity second order derivatives lie in $\cC_\ell^\alpha$. We remark that this fact is (almost) a consequence of the Schauder estimate contained in Theorem~\ref{thm:variable} (as a matter of fact, it is closer to the local version of this result, see Theorem~\ref{thm:variable-loc}). However we do not need here to control the $L^\infty$ norm of $g$. The proof of the following lemma illustrates the hypoelliptic structure of the H\"older spaces $\cC_\ell^\beta$ and is consequently of independent interest.


\begin{lemm}[Hypoelliptic H\"older estimate]\label{l:hypo}
Let $Q$ be an arbitrary cylinder and $g \in \cC_\ell^{2+\alpha} (\R^{2d+1})$ then
\begin{equation}\label{e:hypo-g}
[g]_{\cC_\ell^{2+\alpha}\left(\R^{2d+1}\right)} \le C \left(\left[\left(\partial_t + v \cdot \nabla_x\right) g\right]_{\cC_\ell^{\alpha}\left(\R^{2d+1}\right)}+ \left[D^2_v g\right]_{\cC_\ell^{\alpha}\left(\R^{2d+1}\right)}\right)
\end{equation}
for some constant $C$ only depending on $\alpha$.
\end{lemm}

\begin{rema}\label{rem:equiv}
By combining~\eqref{e:hypo-g} with~\eqref{eq:interp4-bis} and~\eqref{eq:interp5-bis}, we deduce that
\[
\left[\left(\partial_t + v \cdot \nabla_x\right) g\right]_{\cC_\ell^{\alpha}\left(\R^{2d+1}\right)} + \left[D^2_v g\right]_{\cC_\ell^{\alpha}\left(\R^{2d+1}\right)}
\]
is a semi-norm \emph{equivalent} to $[\cdot]_{\cC_\ell^{2+\alpha}(\R^{2d+1})}$. In order words, although the semi-norm $[\cdot]_{\cC_\ell^{2+\alpha}(\R^{2d+1})}$ is defined by measuring oscillations around higher-order polynomials, in the whole space it can be recovered with the more classical notion of H\"older regularity along the highest-order derivatives. We do not know if this equivalence is true in a domain, as in
\[
\left[\left(\partial_t + v \cdot \nabla_x\right) g\right]_{\cC_\ell^{\alpha}(\cQ)} + \left[D^2_v g\right]_{\cC_\ell ^\alpha (\cQ)} \sim [g]_{\cC_\ell^{2+\alpha}(\cQ)}\,.
\]
\end{rema}

\begin{proof}
It is convenient to write $Y$ for $(\partial_t + v \cdot
\nabla_x)$. Moreover, since the domain is the whole space $\R^{2d+1}$, we do not specify the domain in function spaces. Let $z_0 \in \R^{2d+1}$ and $r>0$. We consider in particular $h =(1,u,v) \in Q_1$ and we have $z_0 \circ S_r (h) \in Q_r(z_0)$.

Proving~\eqref{e:hypo-g} for variations along free streaming is easy: for all $z_0 \in \R^{2d+1}$ and $r>0$, we have
\[
g \left(z_0 \circ \left(r^2,0,0\right)\right) - g(z_0) - Yg (z_0) r^2 = r^2 \int_0^1 \left[Y g\left(z_0 \circ \left(r^2 \theta,0,0\right)\right) -Yg(z_0) \right] \dd \theta
\]
which implies
\begin{equation}\label{eq:taylor-fs}
\left|g \left(z_0 \circ \left(r^2,0,0\right)\right) - g(z_0)-Yg (z_0) r^2\right| \le [Yg]_{\cC_\ell^\alpha } r^{2+\alpha}.
\end{equation}

The proof of~\eqref{e:hypo-g} for variations along the $v$ variable is also straightforward: for all $z_0 \in \R^{2d+1}$ and $r>0$,
\begin{equation}\label{eq:taylor-v}
\left|g(z_0 \circ (0,0,r v)) -g(z_0) - \nabla_v g(z_0)
\cdot (rv) - \frac{r^2}2 v^TD^2_v g(z_0)v \right|
\le \left[D^2_v g\right]_{\cC_\ell^{\alpha}} r^{2+\alpha}.
\end{equation}

We then prove~\eqref{e:hypo-g} but only for variations along $x$: we prove that for all $z_0 \in \R^{2d+1}$ and $u \in B_1$ and $r>0$,
\begin{equation}\label{e:holder-x}
\left|g\left(z_0 \circ \left(0,r^3u,0\right)\right) -g(z_0) \right|
\le C \left(\left[(\partial_t + v \cdot \nabla_x) g \right]_{\cC_\ell^{\alpha}} + \left[D_v^2 g\right]_{\cC_\ell^{\alpha}}\right) r^{2+\alpha}
\end{equation}
for some constant $C$ only depending on $\alpha$. This is where a hypoelliptic argument is used: $\nabla_x$ is realized as the Lie bracket of $\nabla_v$ and $Y$, along trajectories.
\[
\xymatrixcolsep{7pc} \xymatrixrowsep{4pc} \xymatrix{ z_1
\ar[r]^{\text{{\tiny forward along $v$}}} & z_2
\ar[r]^{\text{{\tiny backward along transport}}}& z_3 \ar[d]^{\text{\tiny backward along $v$}}\\
& z_0 \ar[lu]^{\text{{\tiny forward along $x$\qquad}}}& z_4
\ar[l]^{\text{{\tiny forward along transport}}} }
\]
Let $z_1$ denote $z_0 \circ (0,r^3u,0)$ and $z_2 = z_1 \circ (0,0,ru)$ and $z_3 = z_2 \circ (-r^2,0,0)$ and $z_4 = z_3 \circ (0,0,-ru)$. Remark that $z_0 = z_4 \circ (r^2,0,0)$. Notice that all points remain in $Q_r (z_0)$. We now write
\begin{multline*}
g(z_0) - g(z_1) =  \left[g\left(z_4 \circ \left(r^2,0,0\right)\right) -g (z_4) \right] +
\left[g(z_3 \circ (0,0,-ru)) -g(z_3) \right]
\\  + \left[g\left(z_2 \circ \left(-r^2,0,0\right)\right) -g(z_2) \right] -
\left[g(z_2\circ (0,0,-ru))-g(z_2)\right].
\end{multline*}

\goodbreak
We now use~\eqref{eq:taylor-fs} with $z_0$ replaced successively by $z_2$ and $z_4$ and~\eqref{eq:taylor-v} with $z_0$ replaced successively by $z_3$ and $z_2$, and we get, after summing the four resulting inequalities,
\begin{equation}\label{e:almost}
\begin{aligned}
\left|g(z_0) -g(z_1)\right| \le & 2 \left([Yg]_{\cC^\alpha_\ell} + \left[D^2_v g\right]_{\cC^\alpha_\ell}\right) r^{2+\alpha} + \left| Yg (z_4)-Yg(z_2)\right| r^2 \\
&\quad+ r \left| \left(\nabla_vg (z_2)-\nabla_v g (z_3)\right) \cdot u \right| + \frac{r^2} 2 \left| D_v^2 g (z_3)-D_v^2 g(z_2)\right| \\
\left|g(z_0) -g(z_1)\right| \le & 3\left([Yg]_{\cC^\alpha_\ell} + \left[D^2_v g\right]_{\cC^\alpha_\ell}\right) r^{2+\alpha}+ \left| \left(\nabla_vg (z_2)-\nabla_v g (z_3)\right) \cdot u \right|
\end{aligned}
\end{equation}
We now estimate $| (\nabla_vg (z_2)-\nabla_v g (z_3)) \cdot u |$ in terms of
\[
[g]_{\cC^{2+\alpha}_{\ell,\,x} } := \sup_{z\,\in\,\R^{2d+1}, \,u\,\in\,B_1} \frac{|g(z \circ (0,0,ru)) - g(z)|}{r^{2+\alpha}}.
\]
We are going to prove that for all $R>0$,
\begin{multline}\label{e:dvx}
\left| \left(\nabla_vg (z_2)-\nabla_v g (z_3)\right) \cdot u \right|\\
\le \left[D^2_v g\right]_{\cC^\alpha_\ell} \left(2R^{1+\alpha}+Rr^\alpha\right) + 3 [Yg]_{\cC^\alpha_\ell} R^{-1} r^{2+\alpha} + [g]_{\cC^{2+\alpha}_{\ell,\,x}} R^{-1} \left(r^2 R\right)^{\frac{2+\alpha}3}.
\end{multline}
In order to derive this estimate, we approximate $\nabla_v g (z) \cdot u$ by the finite difference $\sigma^3_{R,\,u}[g](z)$ we already used above. Recall that
\[
\sigma^3_{R,\,u} [g](z) = \frac{|g(z \circ (0,0,Ru)) - g(z)|}{R}.
\]
Remark that~\eqref{eq:taylor-v} implies
\begin{equation}\label{e:1}
\left|\sigma^3_{R,\,u} [g] (z) - \nabla_v g (z) \cdot u - \frac{R}2 D^2_v g (z) u \cdot u \right| \le \left[D^2_v g\right]_{\cC^\alpha_\ell} R^{1+\alpha}.
\end{equation}
Moreover, recalling that $z_3 = z_2 \circ (-r^2,0,0)$ and using~\eqref{eq:taylor-fs} twice,
\begin{multline} \label{e:2}
 R \left|\sigma^3_{R,\,u} [g](z_2) - \sigma^3_{R,\,u} [g](z_3) \right|
 = \left| g(z_3) - g(z_2) - Yg (z_2) \left(-r^2\right) \right| \\
\begin{aligned}
&\quad + \left|g\left(z_2 \circ (0,0,Ru) \circ (-r^2,0,0)\right) - g\left(z_2 \circ (0,0,Ru)\right)\right.\\
&\left.\quad - Yg \left(z_2 \circ (0,0,Ru)\right)(-r^2) \right| + r^2 \left|Yg (z_2) - Yg \left(z_2 \circ (0,0,Ru)\right)\right| \\
&\quad + \left| g \left(z_2 \circ (0,0,Ru) \circ (-r^2, 0,0)\right) - g\left(z_3 \circ (0,0,Ru)\right)\right|\\
& \le 3 [Yg]_{\cC^\alpha_\ell } r^{2+\alpha} \\
&\quad + \left| g\left(t_2 -r^2,x_2 -r^2 (v_2 + Ru), v_2 + Ru\right) - g\left(t_2-r^2, x_2 -r^2 v_2, v_2 + Ru\right) \right|\\
& \le 3 [Yg]_{\cC^\alpha_\ell} r^{2+\alpha} + [g]_{\cC^{2+\alpha}_{\ell,\,x}} (r^2 R)^{\frac{2+\alpha}3}.
\end{aligned}
\end{multline}
Combining~\eqref{e:1} and~\eqref{e:2} yields~\eqref{e:dvx}.

We now combine~\eqref{e:almost} and~\eqref{e:dvx} and get
\begin{multline*}
r^{-2-\alpha}\left|g(z_0)-g(z_1)\right| \le 3\left([Yg]_{\cC^\alpha_\ell} + \left[D^2_v g\right]_{\cC^\alpha_\ell}\right) \\
+ \left[D^2_v g\right]_{\cC^\alpha_\ell} \left(2 (R/r)^{1+\alpha} + (R/r)\right)
+ 3 [Yg]_{\cC^\alpha_\ell} (r/R) + [g]_{\cC^{2+\alpha}_{\ell,\,x}} (r/R)^{1-\alpha}.
\end{multline*}
We now pick $R = \delta^{-1} r$ with $\delta^{1-\alpha} = 1/2$ and get~\eqref{e:holder-x} for some constant only depending on $\alpha$. This concludes the proof.
\end{proof}


\section{The Schauder estimate}\label{sec:schauder}

This section is devoted to the proof of the Schauder estimate (Theorem~\ref{thm:variable}). The proof proceeds in mainly two steps: in the first step, the matrix $A$ is constant. The estimate is first obtained when $A$ is the identity matrix (Theorem~\ref{thm:constant}), then for an arbitrary constant diffusion matrix (Corollary~\ref{cor:constant}). Then the estimate is established for variable coefficient by the procedure of freezing coefficients thanks to interpolation inequalities.


\subsection{A gradient bound for the Kolmogorov equation}

We first establish a gradient bound for solutions of~\eqref{eq:linear} when $A$ is the identity matrix and when there is no lower order terms ($b=0$, $c=0$). The equation is then reduced to~\eqref{eq:kolmogorov}. Recall that $Q_1= (-1,0] \times B_1 \times B_1$ is the cylinder of radius $1$.

\begin{prop}[Gradient bound]\label{prop:grad}
Given $S \in W^{1,\infty} (Q_1)$, consider $g$ solution to~\eqref{eq:kolmogorov} in $Q_1$, i.e.
\[
(\partial_t + v \cdot \nabla_x)g = \Delta_v g + S.
\]
Then
\begin{multline*}
\forall \, i =1,\,\dots,\,d, \quad \left|\partial_{x_i} g(0,0,0)\right| + |\partial_{v_i} g(0,0,0)| \\
\lesssim_{d}\|g\|_{L^\infty(Q_1)}+\|S\|_{L^\infty(Q_1)} +\|\partial_{x_i} S\|_{L^\infty(Q_1)}+\|\partial_{v_i} S\|_{L^\infty(Q_1)}.
\end{multline*}
\end{prop}

\begin{rema}
See also~\cite{baudoin2017,guillin2012} for related gradient estimates.
\end{rema}


\begin{proof}
We use Bernstein's method as Krylov does in~\cite{krylov1996} in the elliptic-parabolic case, combined with methods from hypocoercivity theory (see for instance~\cite{villani2009memoir}) in order to control the full $(x,v)$-gradient of the solution: see the construction of the quadratic form $w$ in $\partial_{x_i} g$ and $\partial_{v_i} g$ below.

Denote the Kolmogorov operator $\cL_K g := \partial_t g + v \cdot \nabla_x g - \Delta_v g$ and compute the following defaults of distributivity of the operator (reminiscent of the so-called $\Gamma$-calculus~\cite{MR889476})
\begin{equation}\label{eq:square}
\cL_K(g_1g_2) = g_1 \cL_K g_2 + g_2 \cL_K g_1 - 2 \nabla_v g_1\,\cdot\,\nabla_v g_2,\quad \cL_K (g^2) = 2 g \cL_K g - 2 \left| \nabla_v g \right|^2.
\end{equation}


Consider a cut-off function $0 \le \zeta \in C^\infty$ such that $\zeta^{1/2} \in C^\infty$, with support in $(-1,0] \times B_1 \times B_1$ and such that $\zeta(0,0,0)=1$. In order to estimate the gradient of $g$ in $x$ and $v$ at the origin, it is enough to find $\nu_0, \nu_1>0$ and $0 < \frka \le \frkb$ and $0 < \frkc < \frka \frkb$ such that, for any $i \in \{1,\,\dots,\,d\}$,
\[
w = \nu_0 g^2 - \nu_1 t + \left[\frka^2 \zeta^4 (\partial_{x_i}g)^2 + \frkc \zeta^3 (\partial_{x_i} g)(\partial_{v_i}g) + \frkb^2 \zeta^2 (\partial_{v_i}g)^2\right]
\]
satisfies
\[
- \cL_K w \ge 0.
\]
Indeed, the maximum principle for parabolic equations then implies that $\sup_{Q_1} w \linebreak= \sup_{\partial_p Q_1} w$ where $\partial_p Q_1 = \{-1\} \times B_1 \times B_1 \cup [-1,0] \times S_1 \times S_1$ (parabolic boundary). Since $\zeta \equiv 0$ in $\partial_p Q_1$ and $\zeta (0,0,0) =1$, we get
\begin{multline}\label{eq:cs}
\frka^2 \left[\left(\partial_{x_i}g\right)^2 + (\partial_{v_i}g)^2
\right](0,0,0)\\
\begin{aligned}
 &\le \left[\frka^2 (\partial_{x_i}g)^2 + \frkb^2 (\partial_{v_i}g)^2
\right](0,0,0)+ 2 \nu_0 g^2(0,0,0) \\
& \le 2 \left[\frka^2 (\partial_{x_i}g)^2 +
\frkc (\partial_{x_i} g \partial_{v_i}g) + \frkb^2 (\partial_{v_i}g)^2
\right](0,0,0) + 2 \nu_0 g^2(0,0,0) \\
&\le 2 \left[\frka^2 \zeta^4(\partial_{x_i}g)^2 + \frkc \zeta^3(\partial_{x_i} g)(\partial_{v_i}g) + \frkb^2 \zeta^2(\partial_{v_i}g)^2 \right](0,0,0) + 2 \nu_0 g^2(0,0,0) \\
& \le 2 w(0,0,0) \le 2 \left(\nu_0 \sup_{Q_1} g^2 + \nu_1\right).
\end{aligned}
\end{multline}

\begin{enumerate}\romanenumi
%(I)
\item \label{proofprop3.1.i} Observe first that $-\cL_K (-\nu_1t) = \nu_1$.
%\smallskip

%(II)
\item \label{proofprop3.1.ii} Compute second $-\cL_K (\nu_0 g^2)$ using~\eqref{eq:square}
\begin{equation}\label{eq:mu}
- \cL_K(\nu_0 g^2) = 2\nu_0 |\nabla_v g|^2 - 2 S \nu_0 g.
\end{equation}
%\smallskip

%(III)
\item \label{proofprop3.1.iii}Compute third $- \cL_K(\zeta^4 (\partial_{x_i}g)^2)$ using that $\mathcal L_K(\partial_{x_i}g)=\partial_{x_i}S$ and~\eqref{eq:square}

\begin{equation}\label{eq:zeta}
\begin{aligned}
- \cL_K\left(\zeta^4(\partial_{x_i}g)^2\right) & = 2 \zeta^4 \left|\nabla_v\partial_{x_i}g\right|^2 - (\partial_{x_i}g)^2 \cL_K \left(\zeta^4\right)
\\
& \quad + 2 \nabla_v (\zeta^4) \cdot
\nabla_v (\partial_{x_i} g)^2 - 2 \zeta^4 \partial_{x_i}g
\partial_{x_i}S \\
& \ge \zeta^4 \left|\nabla_v
\partial_{x_i}g\right|^2 + (\partial_{x_i}g)^2
\left[- \cL_K (\zeta^4) - 4 \zeta^{-4} \left|\nabla_v (\zeta^4)\right|^2 \right] \\
& \quad - \zeta^3 (\partial_{x_i}g)^2 - \zeta^5 (\partial_{x_i}S)^2.
\end{aligned}
\end{equation}
%\smallskip

%(IV)
\item \label{proofprop3.1.iv}Compute fourth, for some $\epsilon_1>0$, the term $- \cL_K(\zeta^2(\partial_{v_i}g)^2)$ using
\[
\cL_K (\partial_{v_i} g) = \partial_{v_i}S -
\partial_{x_i}g
\]
and~\eqref{eq:square}:
\begin{multline}\label{eq:barbarzeta}
- \cL_K\left(\zeta^2 (\partial_{v_i}g)^2\right)
 = 2 \zeta^2 \left|\nabla_v \partial_{v_i} g\right|^2 + 2 \nabla_v (\zeta^2) \cdot \nabla_v (\partial_{v_i} g)^2 \\
\begin{aligned}
& \mkern50mu- (\partial_{v_i} g)^2 \cL_K (\zeta^2) + 2 \zeta^2
\partial_{v_i} g \partial_{x_i} g -2 \zeta^2 \partial_{v_i}g \partial_{v_i} S \\
& \ge \zeta^2 \left|\nabla_v \partial_{v_i} g\right|^2 + (\partial_{v_i} g)^2 \left[- 1 -
\epsilon_1 ^{-1} \zeta -
\cL_K (\zeta^2) - 2 \zeta^{-2} \left|\nabla_v
\zeta\right|^2\right] \\
& \mkern50mu - \epsilon_1 \zeta^3 (\partial_{x_i}g)^2 - \zeta^4 (\partial_{v_i} S)^2.
\end{aligned}
\end{multline}
%\smallskip

%(V)
\item \label{proofprop3.1.v}Compute fifth, for some $\epsilon_2>0$, the term $-\cL_K [\zeta^3 (\partial_{x_i}g) (\partial_{v_i} g)]$, with the intermediate step
\[
- \cL_K \left[(\partial_{x_i} g)(\partial_{v_i} g)\right] = (\partial_{x_i} g)^2 + 2 \nabla_v \partial_{x_i} g \cdot \nabla_v \partial_{v_i} g -\partial_{x_i}g \partial_{v_i}S - \partial_{v_i} g \partial_{x_i}S,
\]
\begin{multline}\label{eq:barzeta}
- \cL_K \left[\zeta^3 (\partial_{x_i} g)\left(\partial_{v_i} g\right)\right]\\
\begin{aligned}
= &\zeta^3 \left[(\partial_{x_i} g)^2 + 2 \nabla_v \partial_{x_i} g \cdot \nabla_v \partial_{v_i} g
\right] - \left(\partial_{x_i}g\right)\left(\partial_{v_i} g\right) \cL_K (\zeta^3) \\
&+ 2 \nabla_v (\zeta^3) \cdot \nabla_v \left[(\partial_{x_i} g)(\partial_{v_i}g)\right] - \zeta^3 \partial_{x_i}g
\partial_{v_i}S - \zeta^3 \partial_{v_i}g \partial_{x_i}S \\
\ge & \frac12 \zeta^3 \left(\partial_{x_i} g\right)^2 - \epsilon_2 \zeta^4 \left|\nabla_v
\partial_{x_i}g\right|^2 - \epsilon_2^{-1} \zeta^2 \left|\nabla_v \partial_{v_i} g \right|^2
-(\partial_{x_i}g)(\partial_{v_i} g) \cL_K (\zeta^3) \\
& + 2 \nabla_v (\zeta^3) \cdot \nabla_v \left[(\partial_{x_i} g)\right] (\partial_{v_i}g) + 2 \nabla_v (\zeta^3) \cdot \nabla_v \left[(\partial_{v_i}g)\right] \left(\partial_{x_i} g\right) \\
& - \frac12 \zeta^3 \left(\partial_{v_i}S\right)^2 - \frac12 \zeta^3 \left(\partial_{v_i}g\right)^2 - \frac12 \zeta^3 \left(\partial_{x_i}S\right)^2.
\end{aligned}
\end{multline}
\end{enumerate}


To clean a little the calculations, observe that
\begin{enumerate}
\item \label{proofprop3.1.1}error terms involving $S$ are controlled by choosing $\nu_1$ larger enough, i.e. larger than a multiple of $\|S\|_{L^\infty(Q_1)} +\|\partial_{x_i} S\|_{L^\infty(Q_1)}+\|\partial_{v_i} S\|_{L^\infty(Q_1)}$,
\item \label{proofprop3.1.2}terms involving $|\nabla_v g|^2$ or $\nabla_v g \cdot \nabla_x g$ are controlled by choosing $\nu_0$ large enough thanks to~\eqref{eq:mu}, with $g$ controlled by its sup norm,
\item \label{proofprop3.1.3}Equation~\eqref{eq:mu} is ``free'' (i.e. not involved in any constant dependency) as well as equation~\eqref{eq:barbarzeta} by choosing $\epsilon_1$ small enough so that the term $-\epsilon_1 \zeta^3 (\partial_{x_i} g)^3$ is compensated by the positive term in~\eqref{eq:barzeta},
\item \label{proofprop3.1.4}Equation~\eqref{eq:zeta} has an error term of the form $-O(1) \zeta^3 (\partial_{x_i} g)^2$ that is compensated by the positive term in~\eqref{eq:barzeta} again: we use $|- \cL_K (\zeta^4) - 2 \zeta^{-4} |\nabla_v (\zeta^4)|^2 |\linebreak
\lesssim \zeta^3$ (remember that $\zeta^{\frac12} \in C^\infty$).
\item \label{proofprop3.1.5}In the last equation~\eqref{eq:barzeta} we also split
\begin{multline*}
- (\partial_{x_i}g)(\partial_{v_i} g) \cL_K (\zeta^3)
\lesssim \frac18 (\partial_{x_i} g)^2 + O(1) (\partial_{v_i} g)^2,\\
\begin{aligned}
& 2 \nabla_v (\zeta^3) \cdot \nabla_v [(\partial_{x_i} g)] (\partial_{v_i}g) \lesssim \epsilon_3 \zeta^4 |\nabla_v [(
\partial_{x_i} g)]|^2 + \epsilon_3 ^{-1} O(1) (\partial_{v_i}g) ^2 \\
& 2 \nabla_v (\zeta^3) \cdot \nabla_v [(\partial_{v_i}g)] (
\partial_{x_i} g) \lesssim O(1) \zeta^2 |\nabla_v [(\partial_{v_i}g)] (\partial_{x_i} g)|^2 + \frac{1}{4}
\zeta^3(\partial_{x_i} g)^2,
\end{aligned}
\end{multline*}
where we have used again $\zeta^{\frac12} \in C^\infty$, in the form $|\nabla \zeta| \lesssim \zeta^{1/2}$.
\end{enumerate}
%\medskip

These considerations result in the following calculations:
\begin{multline*}
- \cL_K w \ge 2 \nu_0 \left|\nabla_v g\right|^2 + \nu_1 + \frka^2 \zeta^4 \left|\nabla_v
\partial_{x_i}g\right|^2 - \frka^2 O(1) \zeta^3 \left(\partial_{x_i}g\right)^2 - \frka^2 O(1) (\partial_{x_i}S)^2 \\
\begin{aligned}
& + \frkb^2 \zeta^2 \left|\nabla_v \partial_{v_i} g\right|^2 - \frkb^2 O(1) \left(\partial_{v_i} g\right)^2 - \frkb^2 \epsilon_1 \zeta^3 \left(\partial_{x_i}g\right)^2 - \frkb^2 O(1) \left(\partial_{v_i} S\right)^2 \\
& + \frac{\frkc}4 \zeta^3 \left(\partial_{x_i} g\right)^2 - \frkc \left(\epsilon_2 +\epsilon_3\right) \zeta^4 \left|\nabla_v\partial_{x_i}g\right|^2 - \frkc O(1) \zeta^2 \left|\nabla_v \partial_{v_i} g \right|^2 \\
& - \frkc O(1) \left(\partial_{x_i,\,v_i}S\right)^2 - \frkc O(1) \left(\partial_{v_i}g\right)^2.
\end{aligned}
\end{multline*}
We finally choose
\begin{enumerate}
\item \label{proofprop3.1.2.1}$\frka=1$,
\item \label{proofprop3.1.2.2}$\frkc$ large enough so that the first term in the third line controls the fourth term in the first line,
\item \label{proofprop3.1.2.3}$\epsilon_2$ and $\epsilon_3$ small enough so that the second term of the third line is controlled by the second term in the right hand side of the first line,
\item \label{proofprop3.1.2.4}$\frkb$ large enough so that the third term in the fourth line is controlled by the first term in the third line and $\frka \frkb > \frkc$ so that the quadratic form is strictly positive,
\item \label{proofprop3.1.2.5}$\epsilon_1$ small enough so that the third term in the third line is controlled by the first term in the third line,
\item \label{proofprop3.1.2.6}finally $\nu_0$ and $\nu_1$ large enough to control all the $v$-derivatives of $g$ and all the derivatives of $S$. Notice that
\[
\nu_1 = O(1)\left(\| g \|_{L^\infty(Q_1)} + \|S\|_{L^\infty(Q_1)}^2 + \|\partial_{x_i} S\|_{L^\infty(Q_1)}^2+ \|\partial_{v_i} S\|_{L^\infty(Q_1)}^2\right).
\]
\end{enumerate}
The conditions on the coefficients are:
\begin{align*}
\epsilon_2, \epsilon_3 <<1, \quad \frkc >> 1, \quad
\frkb^2 >> \frkc, \quad
\epsilon_1 << \frac{\frkc}{\frkb^2}, \quad
\nu_0, \nu_1 >> \frkb^2, \frkc
\end{align*}
which is compatible with all requirements and has solutions. This proves that $-\cL_K \omega \ge 0$ and the desired inequality is thus obtained from~\eqref{eq:cs} and the choice of $\nu_1$ we made above, which concludes the proof of Proposition~\ref{prop:grad}.
\end{proof}

A direct consequence of the gradient estimate from Proposition~\ref{prop:grad} is a bound of derivatives of arbitrary order in any cylinder of radius $r$. We recall that $Q_r\linebreak= (-r^2,0] \times B_{r^3} \times B_r$.

\begin{coro}[Bounds on arbitrary derivatives around the
origin]\label{cor:deriv}
Given $k \in \N$ and $r>0$, there exists a constant $C$ depending on dimension $d$ and $k$ such that any solution of~\eqref{eq:kolmogorov} in $Q_r$ with zero source term $S \equiv 0$, i.e. $(\partial_t + v \cdot \nabla_x) g = \Delta_v g$ in $Q_r$, satisfies for all $n \in \N$ and $\alpha,\beta \in \N^d$ with $|\beta|=k$,
\[
\left|\partial_t^n D^\alpha_x D^\beta_v g(0,0,0)\right| \le \frac{C
\|g\|_{L^\infty(Q_r)}}{r^{2n +3|\alpha|+|\beta|}}.
\]
\end{coro}


\begin{rema}
Remark that $2n+3|\alpha|+|\beta|$ is the kinetic degree of the polynomial associated with $\partial_t^n D^\alpha_x D^\beta_v$.
\end{rema}


\begin{proof}
We reduce to the case $r=1$ by rescaling: the function $g_r (t,x,v) \linebreak= g(r^2 t, r^3 x, r v)$ is a solution of~\eqref{eq:kolmogorov} in $Q_1$. If the result is true for $r=1$, then we get the desired estimate for arbitrary $r$'s. We then first treat the case $n=0$ and argue by induction on $|\beta|$. Proposition~\ref{prop:grad} yields the result for $|\beta|\le 1$ since $D^\alpha_x g$ solves~\eqref{eq:kolmogorov} with $S \equiv 0$ for an arbitrary multi-index $\alpha$. Assuming the result true for $n=0$, any $\alpha$ and $|\beta| \le k$, remark that $D^\alpha_x D_v ^\beta g$ solves~\eqref{eq:kolmogorov} with $S =\sum_{i=1,\,\dots,\,d \, | \, \beta_i \ge 1} D^{\alpha+\delta_i} _x D^{\beta-\delta_i}_ v g$. Consider $\beta \in \N^d$ with $|\beta|=k+1$; the previous step yields controls of $\partial_{x_j} S$ and $\partial_{v_j} S$ for the source term $S$ in the equation for $D^\alpha_x D_v ^\beta g$. Proposition~\ref{prop:grad} then gives the control $\partial_{x_i,v_i} D^\alpha_x D_v ^\beta g(0,0,0)$ which completes the induction. We finally get the result for an arbitrary $n \ge 1$ by remarking that the equation allows us to control any time derivatives by space and velocity derivatives.
\end{proof}


\subsection{Proof of Schauder estimates}

With such bounds on derivatives at hand, we can turn to the proof of the Schauder estimate for the Kolmogorov equation, that is to say for equation~\eqref{eq:linear} with $A$ replaced with the identity matrix and with no lower order terms ($b=0$, $c=0$), see~\eqref{eq:kolmogorov}. A change of variables will then yield the result for any constant matrix $A$ satisfying the ellipticity condition~\eqref{eq:elliptic}. Finally we shall classically approximate the coefficients locally by constants to treat the general case.

\subsubsection{The core estimate}

The proof of the Schauder estimate for the Kolmogorov equation follows the argument proposed by Safonov~\cite{safonov1984}, as explained in~\cite{krylov1996}.
\begin{theo}[The Schauder estimate for the Kolmogorov
equation]\label{thm:constant}
Given $\alpha \in (0,1)$ and $S \in \cC_\ell^\alpha$, let $g \in \cC_\ell^{2+\alpha}$ be a solution to~\eqref{eq:kolmogorov}, i.e. $(\partial_t + v \cdot \nabla_x) g = \Delta_v g + S$ in $\R \times \R^d \times \R^d$. Then
\[
[g]_{\cC_\ell^{2+\alpha}} \lesssim_{d,\,\alpha} [S]_{\cC_\ell^\alpha}.
\]
\end{theo}

\begin{rema}
This theorem is the hypoelliptic counterpart of~\cite[Theorem~8.6.1~\& Lemma~8.7.1]{krylov1996}.
\end{rema}

\begin{proof}
We first reduce to the case where $g \in C^\infty_c (\R \times \R^d \times \R^d)$ by mollification and truncation (as for instance in the proof of~\cite[Lemma~8.7.1, p.~122]{krylov1996}). We then reduce to the base point $z_0=0$ by considering the change of unknown $g^\sharp(z) := g(z_0 \circ z)$ (we however keep on calling the unknown $g$). Given $r>0$ and $K \ge 1$ to be chosen later, consider $Q_{(K+1)r}$ and a cut-off function $\zeta \in C^\infty_c$ such that $\zeta \equiv 1$ in $Q_{(K+1)r}$. Recall again the Kolmogorov operator $\cL_K := \partial_t + v \cdot \nabla_x - \Delta_v$, and define $\bar S := \cL_K (\zeta \cT_{0}[g])$ with the Taylor polynomial $\cT_{z_0}[g]$ of $g$ at $(0,0,0)$, defined in~\eqref{e:taylor}.

Decompose in $Q_{(K+1)r}$ (where $\zeta =1$):
\begin{multline*}
g - \cT_{0}[g] = g - \zeta
\cT_{0}[g] = \cG \starkin (S -\bar S) = h_1 + h_2
\\
\text{ with } \quad
\begin{cases} h_1 := \cG \starkin \left[(S- \bar S)
\un_{Q_{(K+1)r}} \right]
 \\[2mm] h_2 := \cG \starkin \left[(S-\bar S) \un_{Q^c_{(K+1)r}} \right]
\end{cases},
\end{multline*}
where $Q^c_{(K+1)r} = \R \times \R^d \times \R^d \setminus Q_{(K+1)r}$ and $\cG$ is the Green function studied in Proposition~\ref{prop:pregreen}.

Observe that $h_1$ is the solution to $\partial_t h_1 + v \cdot \nabla_x h_1 -\Delta_v h_1 = S-S(0,0,0)$ since
\begin{align*}
\bar S(z) &= \cL_K(\zeta \cT_0[g])(z) = \cL_K(\cT_0[g])(z) \\
&= \left(\partial_t + v \cdot \nabla_x g\right)(0,0,0) - \Delta_v g(0,0,0) = \text{cst} = S(0,0,0)
\end{align*}
for $z \in Q_{(K+1)r}$, and $h_2$ is the solution to $\partial_t h_2 + v \cdot \nabla_x h_2- \Delta_v h_2 = 0$ in $Q_{(K+1)r}$.

We next estimate
\begin{equation}\label{eq:start}
\left\| g - \cT_{0}[g] - \cT_{0}[h_2] \right\|_{L^\infty(Q_r)} \le \left\| h_2 - \cT_{0}[h_2] \right\|_{L^\infty(Q_r)} + \| h_1\|_{L^\infty(Q_r)}.
\end{equation}
Using Proposition~\ref{prop:pregreen} and $S \in \cC^\alpha_\ell$, we get
\begin{equation}\label{eq:g}
\| h_1 \|_{L^\infty(Q_r)} \lesssim (K+1)^{2+\alpha} r^{2+\alpha} [S]_{\cC_\ell^{\alpha}\left(Q_{(K+1)r}\right)}.
\end{equation}
Now for $z = (r^2 t, r^3 x, rv) \in Q_r$ with $(t,x,v) \in Q_1$. There exists $\theta_1, \theta_2, \theta_3 \in (0,1)$ such that
\begin{align*}
h_2(z) = & h_2\left(r^2 t, r^3 x, rv\right) \\
= & h_2\left(r^2 t, 0, rv\right) + \left(r^3 x\right) \cdot \nabla_x h_2 \left(r^2t,r^3 \theta_1 x,r v\right) \\
= & h_2(0, 0, rv) + \left(r^3 x\right) \cdot \nabla_x h_2 \left(r^2t,r^3 \theta_1 x,r v\right) + r^2t \partial_t h_2 \left(\theta_2 r^2 t, 0,r v\right) \\
= & h_2(0,0,0) + \left(r^3 x\right) \cdot \nabla_x h_2 \left(r^2t,r^3 \theta_1 x,r v\right) + r^2t \partial_t h_2 \left(\theta_2 r^2 t, 0,r v\right) \\
& \quad + \nabla_v h_2(0,0,0) \cdot (rv) + \frac12 (rv) ^T \cdot D^2_v h_2(0,0, r \theta_3 v) \cdot (rv).
\end{align*}
As a consequence
\begin{multline*}
\left\| h_2 - \cT_{0}[h_2]\right\|_{L^\infty(Q_r)} \le r^2 \left[r \left\|\nabla_x h_2\right\|_{L^\infty\left(Q_{(K+1)r}\right)} + r^2 \left\|\partial^2_t h_2\right\|_{L^\infty\left(Q_{(K+1)r}\right)}\right. \\
+\left. r \left\|\partial_t \nabla_v h_2 \right\|_{L^\infty\left(Q_{(K+1)r}\right)} + r \left\| D^3_v h_2 \right\|_{L^\infty\left(Q_{(K+1)r}\right)} \right].
\end{multline*}

We remark that $h_2$ satisfies~\eqref{eq:kolmogorov} with $S\equiv 0$ in $Q_{(K+1)r}$. We thus can apply Corollary~\ref{cor:deriv} and get
\begin{align*}
\left\|h_2 - \cT_{0}[h_2]\right\|_{L^\infty(Q_r)} & \lesssim r^2 \|h_2\|_{L^\infty\left(Q_{(K+1)r}\right)}  \\
& \;\:\times \left[\frac{r}{((K+1)r)^3} + \frac{r^2}{((K+1)r)^4}+\frac{r}{((K+1)r)^3}+\frac{r}{((K+1)r)^3}\right]
\\
& \lesssim (K+1)^{-3} \|h_2\|_{L^\infty\left(Q_{(K+1)r}\right)}.
\end{align*}
Since $g-\cT_{0} g = h_1 +h_2$, we can estimate $\|h_2\|_{L^\infty(Q_{(K+1)r})}$ as follows
\begin{align*}
\|h_2\|_{L^\infty\left(Q_{(K+1)r}\right)} &\le \|h_1\|_{L^\infty\left(Q_{(K+1)r}\right)} + \|g - \cT_{0}[g] \|_{L^\infty\left(Q_{(K+1)r}\right)} \\
& \lesssim (K+1)^{2+\alpha} r^{2+\alpha} \left([S]_{\cC_\ell^{\alpha}\left(Q_{(K+1)r}\right)} + [g]_{\cC_\ell^{2+\alpha}\left(Q_{(K+1)r}\right)} \right)
\end{align*}
(we used~\eqref{eq:g}) and get
\begin{equation}\label{eq:h}
\|h_2 -\cT_{0}[h_2]\|_{L^\infty(Q_r)} \le C
\frac{r^{2+\alpha}}{(K+1)^{1-\alpha}}
\left([S]_{\cC_\ell^{\alpha}\left(Q_{(K+1)r}\right)} + [g]_{\cC_\ell^{2+\alpha}\left(Q_{(K+1)r}\right)}\right).
\end{equation}
Combining~\eqref{eq:start}, \eqref{eq:g} and~\eqref{eq:h}, we get
\begin{multline*}
\left\| g - \cT_{0}[g] - \cT_{0} [h]_2 \right\|_{L^\infty(Q_r)} \\
\begin{aligned}
& \le \left\| h_2 - \cT_{0} h_2 \right\|_{L^\infty(Q_r)} + \| h_1\|_{L^\infty(Q_r)} \\
& \le C (K+1)^{2+\alpha} r^{2+\alpha} [S]_{\cC_\ell^{\alpha}\left(Q_{(K+1)r}\right)} \\
& \quad + C \frac{r^{2+\alpha}}{(K+1)^{1-\alpha}} [S]_{\cC_\ell^{\alpha}\left(Q_{(K+1)r}\right)} + C \frac{r^{2+\alpha}}{(K+1)^{1-\alpha}} [g]_{\cC_\ell^{2+\alpha}\left(Q_{(K+1)r}\right)}
\end{aligned}
\end{multline*}
and by setting $K$ large enough so that $C (K+1)^{-(1-\alpha)} \le \frac12$, it results into
\[
[g]_{\cC_\ell^{2+\alpha}} \le \sup_{z_0,\:r\,>\,0} \frac{ \left\| g -
\cT_{z_0}[g] - \cT_{z_0}[h_2]
\right\|_{L^\infty(Q_r)}}{r^{2+\alpha}} \le C_K [S]_{\cC_\ell^{\alpha}}
\]
where we have used that $\cT_{z_0}[g] + \cT_{z_0}[h_2]$ is a polynomial of kinetic order less or equal to $2$, which concludes the proof. Note the \emph{corrector} $\cT_{z_0}[h_2]$ to the Taylor polynomial $\cT_{z_0}[g]$ used in this argument.
\end{proof}

The general constant coefficients case is reached through a change of variables.
\begin{coro}[Schauder estimates for
constant diffusion coefficients]\label{cor:constant}
Let $\alpha \in (0,1)$ and $S \in \cC_\ell^\alpha$ and let $A:=(a^{i,\,j})$ be a constant real $d \times d$ matrix that satisfies~\eqref{eq:elliptic}. Then for all solution $g \in \cC_\ell^{2+\alpha}$ to
\[
(\partial_t + v \cdot \nabla_x) g = \sum_{1\,\le\:i,\,j\:\le\, d} a^{i,\,j} \partial^2_{v_i v_j}\,g \; \text{ in }\;  \R \times \R^d
\times \R^d
\]
we have,
\[
[g]_{\cC_\ell^{2+\alpha}} \le C [S]_{\cC_\ell^{\alpha}}
\]
where the constant $C$ depends on $d$, $\alpha$, the norm of the (constant) matrix $A$ and $\lambda$ in~\eqref{eq:elliptic}.
\end{coro}
\goodbreak

\begin{rema}
This is the counterpart to~\cite[Theorem~8.9.1, p.~127]{krylov1996}.
\end{rema}


\begin{proof}
Consider the change of variables
\[
\bar g(t,x,v) := g(t,B x,B v)
\]
with $B^2=A^{-1}$. Then we have for $(t,x,v) \in \R \times \R^d \times \R^d$,
\begin{align*}
& (\partial_t + v \cdot \nabla_x) \bar g (t,x,v) - \Delta_v \bar g(t,x,v) \\
&\qquad = (\partial_t + v \cdot \nabla_x) g (t,B x,B v) - \sum_{1\,\le\:i,\,j\:\le\,d} a^{i,\,j} \partial^2_{v_i v_j} g(t,B x, B v) \\
&\qquad = S(t,Bx,Bv) =: \bar S (t,x,v)
\end{align*}
and the result follows from Theorem~\ref{thm:constant}.
\end{proof}


\subsubsection{Proof of Theorem~\ref{thm:variable}}

\begin{proof}[Proof of Theorem~\ref{thm:variable}]
It is enough to treat the case where $b \equiv 0$ and $c \equiv 0$, the general case is then treated by interpolation (using Lemma~\ref{l:interpol}). We consider a constant $\gamma>0$ which will be fixed later, and pick $z_0,z_1 \in \R \times \R^d \times \R^d$ and $r >0$ such that $z_1 \in Q_r(z_2)$ and
\[
[g]_{\cC_\ell^{2+\alpha}} \le 2 \frac{\left| g (z_1) -
\cT_{z_0}[g](z_1)\right|}{r^{2+\alpha}}.
\]
%\smallskip
%\noindent {\bf Case~1: $r \ge \gamma$.}
\goodbreak
\begin{case}{$r \ge \gamma$\label{case1}}
We compute then
\begin{multline*}
[g]_{\cC_\ell^{2+ \alpha}} \\
\begin{aligned}
& \le 2 \left(2 \gamma^{-2-\alpha} \|g\|_{L^\infty}\! + \left\| \left(\partial_t + v \cdot \nabla_x\right) g\right\|_{L^\infty} \gamma^{-\alpha} \!+ \|\nabla_v g \|_{L^\infty} \gamma^{-1-\alpha}\!+ \frac12 \left\| D^2_v g \right\|_{L^\infty} \gamma^{-\alpha}\! \right) \\
& \le \frac14 [g]_{\cC_\ell^{2+\alpha}} + C_1(\gamma) \|g\|_{L^\infty}
\end{aligned}
\end{multline*}
(using again Lemma~\ref{l:interpol}) with $C_1(\gamma) >0$ depending on $\gamma$ and $d$.
\end{case}


%{\bf Case~2: $r \le \gamma$.}
\begin{case}{$r \le \gamma$\label{case2}}

Then $z_1 \in Q_\gamma (z_0)$ and we consider a $C^\infty$ cut-off function $0 \le \xi \le 1$ that is equal to $1$ on $Q_\gamma(z_0)$ and equal to zero outside $Q_{2\gamma}(z_0)$. In particular, $\xi (z_1)\linebreak=\xi(z_0)=1$. We now use Corollary~\ref{cor:constant} to get
\begin{align*}
[g]_{\cC_\ell^{2+ \alpha}} & \le 2 [g \xi]_{\cC_\ell^{2+\alpha}}
\le C \left[(\partial_t + v\cdot \nabla_x) (g\xi) - \sum_{i,\,j} a^{i,\,j}(z_1)
\partial_{v_iv_j} ^2 (g\xi)\right]_{\cC_\ell^{\alpha}} \\
& \le C \left[(\partial_t + v\cdot \nabla_x) (g\xi) - \sum_{i,\,j} a^{i,\,j}(\cdot)
\partial_{v_iv_j} ^2 (g\xi) \right]_{\cC_\ell^{\alpha}} \\
& \qquad + C \left[\sum_{i,\,j} \left(a^{i,\,j}(\cdot) - a^{i,\,j}(z_1)\right)
\partial_{v_iv_j} ^2 (g\xi) \right]_{\cC_\ell^{\alpha}\left(Q_{2 \gamma}(z_2)\right)}
\end{align*}
where we have added the restriction to the support of $\xi$ in the last term. We estimate successively the two terms of the right hand side. On the one hand, recalling that $(\partial_t + v\cdot \nabla_x) g = \sum_{i,j} a^{i,\,j}\partial_{v_iv_j} ^2 g+S$,
\begin{multline*}
 \left[(\partial_t + v\cdot \nabla_x) (g\xi) -\sum_{i,\,j} a^{i,\,j}
\partial^2_{v_i v_j}(g\xi) \right]_{\cC_\ell^{\alpha}} \\
\begin{aligned}
& \le \left[(\partial_t + v\cdot \nabla_x) g - \sum_{i,\,j} a^{i,\,j}
\partial_{v_i v_j} ^2g \right]_{\cC_\ell^{\alpha}} \\
&  \quad + \left[g \left\{ (\partial_t + v\cdot \nabla_x) (\xi) - \sum_{i,\,j} a^{i,\,j}
\partial_{v_i v_j} ^2\xi \right\} \right]_{\cC_\ell^{\alpha}} + \left[\nabla_v g \cdot \nabla_v \xi\right]_{\cC_\ell^{\alpha}} \\
& \lesssim_\gamma [S]_{\cC_\ell^{\alpha}} + \|g\|_{\cC_\ell^{\alpha}} + \| \nabla_v g \|_{\cC_\ell^{\alpha}} \\
& \le \frac14 [g]_{\cC_\ell^{2+\alpha}} + C_2(\gamma) \left([S]_{\cC_\ell^{\alpha}}+\|g\|_{L^\infty}\right)
\end{aligned}
\end{multline*}
(using again Lemma~\ref{l:interpol}) for a constant $C_2(\gamma)>0$ depending on $d$, $\|A\|_{\cC_\ell^{\alpha}}$, $\alpha$. On the other hand,
\begin{align*}
\left[\sum_{i,\,j} \left(a^{i,\,j}(\cdot) - a^{i,\,j}(z_1)\right)
\partial_{v_iv_j} ^2(g\xi) \right]_{\cC_\ell^{\alpha}\left(Q_{2\gamma} (z_2)\right)} & \lesssim \gamma^{\alpha} \left[D^2_v g\right]_{\cC_\ell^{\alpha}} + \left\|D^2_v g \right\|_{L^\infty} \\
& \le C_3 \gamma^\alpha [g]_{\cC_\ell^{2+\alpha}} + C_4(\gamma) \|g\|_{L^\infty}
\end{align*}
using again Lemma~\ref{l:interpol}, for some constants $C_3>0$ and $C_4(\gamma)$ depending $\gamma$. Combining the last three estimates yields finally in the case $r \le \gamma$:
\[
[g]_{\cC_\ell^{2+\alpha}} \le \left(C_3 \gamma^\alpha + \frac14 \right) [g]_{\cC_\ell^{2+\alpha}} + \left(C_2(\gamma) + C_4(\gamma)
\right) \left([S]_{\cC_\ell^\alpha} + \|g\|_{L^\infty}\right).
\]
\end{case}
We now pick $\gamma$ such that $C \gamma^\alpha + \frac14 \le \frac12$ and we thus get in both cases ($r \ge \gamma$ and $r \le \gamma$) that
\[
[g]_{\cC_\ell^{2+\alpha}} \le \frac12 [g]_{\cC_\ell^{2+\alpha}} + C_5(\gamma)
\left([S]_{\cC_\ell^\alpha} + \|g\|_{L^\infty}\right)
\]
for some constant $C_5(\gamma)>0$, which concludes the proof of~\eqref{eq:schauder-main} thanks to~\eqref{eq:interp4-bis} and~\eqref{eq:interp5-bis}.
\end{proof}

\subsection{Localization of the Schauder estimate}\label{sec:local}

\begin{theo}[Local Schauder estimate]\label{thm:variable-loc}
Given $\alpha \in (0,1)$ and $a^{i,j},b^i,c, S \in\cC_\ell^{\alpha}$ satisfying~\eqref{eq:elliptic} for some constant $\lambda >0$, any solution $g \in \cC_\ell^{2+\alpha}$ to~\eqref{eq:linear} satisfies for all $z_0 \in \R \times \R^d \times \R^d$
\begin{align*}
\| g\|_{\cC_\ell^{2+\alpha}\left(Q_1(z_0)\right)}
\le C \left([S]_{\cC_\ell^{\alpha}\left(Q_2(z_0)\right)} + \| g
\|_{L^\infty\left(Q_2(z_0)\right)} \right)
\end{align*}
where the constant $C$ depends on $d$, $\lambda$ and $\alpha$ and $\|a^{i,j}\|_{\cC_\ell^{\alpha}}$, $\|b^i\|_{\cC_\ell^{\alpha}}$, $\|c\|_{\cC_\ell^{\alpha}}$ for $i,j\linebreak =1,\,\dots,\,d$.
\end{theo}
\begin{proof}
We use the strategy of~\cite[Theorem~8.11.1]{krylov1996}. Consider $z_0=0$ without loss of generality and define $R_n := \sum_{j =0}^n 2^{-j}$ for $n \ge 0$. Define a cutoff function $\zeta_n$ that is smooth, one on $Q_{R_n}$ and zero outside $Q_{R_{n+1}}$. It satisfies the controls
\[
\| \zeta_n \|_{\cC_\ell^{\alpha}}, \left\| v \cdot \nabla_x \zeta_n
\right\|_{\cC_\ell^{\alpha}}, \left\| \nabla_v \zeta_n \right\|_{\cC_\ell^{\alpha}}, \left\|
\nabla_v ^2 \zeta_n \right\|_{\cC_\ell^{\alpha}} \lesssim \rho^{-n} \quad
\text{ with } \quad \rho :=2^{2+\alpha}.
\]


Then apply the non-localized estimate of Theorem~\ref{thm:variable} to $\zeta_n g$. In order to do so, it is convenient to consider the differential operator $\cL = (\partial_t + v \cdot \nabla_x) - \sum_{i,j} a^{i,j}
\partial^2_{v_i v_j}\linebreak- \sum_i b^i \partial_{v_i}$, recall $\cL g =S$ and write
\begin{multline*}
A_n := [g]_{\cC_\ell^{2+\alpha}\left(Q_{R_n}\right)} \le [\zeta_n g]_{\cC_\ell^{2+\alpha}}
\lesssim [\cL (\zeta_n g)]_{\cC_\ell^{\alpha}} + \| \zeta_n g \|_{L^\infty} \\
\begin{aligned}
\lesssim & \left[\zeta_n \cL g + g \cL \zeta_n + \left(A \nabla_v g\right) \cdot \nabla_v \zeta_n
\right]_{\cC_\ell^{\alpha}} + \| \zeta_n g \|_{L^\infty} \\
\lesssim & [S]_{\cC_\ell^{\alpha}(Q_2)} + \rho^{-n}
\left(\| g\|_{\cC_\ell^{\alpha}\left(Q_{R_{n+1}}\right)} + \| \nabla_v g\|_{\cC_\ell^{\alpha}\left(Q_{R_{n+1}}\right)}\right) +\|g\|_{L^\infty\left(Q_{R_{n+1}}\right)}
\end{aligned}
\end{multline*}
and use the interpolation inequalities from Lemma~\ref{l:interpol} to get for all $\eps_n >0$,
\[
\mkern60mu\lesssim [S]_{\cC_\ell^{\alpha}(Q_2)} + \rho^{-n} \left(\eps_n A_{n+1} + \eps_n^{-\beta} \rho^{-n} \|g\|_{L^\infty(Q_2)}\right) + \|g\|_{L^\infty\left(Q_{R_{n+1}}\right)}
\]
for some $\beta >0$ (note the additional factor $\rho^{-n}$ due to the dependency of the constant $\tilde{C}$ in the difference of radii in Lemma~\ref{l:interpol}). Choosing next $\eps_n := \eps_0 \rho^n$ for $\eps_0 \in (0,1)$ small enough yields
\begin{align*}
A_n \lesssim [S]_{\cC_\ell^{\alpha}(Q_2)} + \eps_0 A_{n+1} + \eps_0 ^{-\beta} \rho^{-(\beta+2) n} \| g \|_{L^\infty(Q_2)}
\end{align*}
Consider then the geometric sum $\sum_{n\,\ge \,0} \eps_0 ^n A_n$, and calculate
\[
\sum_{n\,\ge\,0} \eps_0^n A_n \lesssim \left(\sum_{n\,\ge\, 0}
\eps_0^n \right) [S]_{\cC_\ell^{\alpha}(Q_{2})} + \sum_{n\, \ge\,0}
\eps_0^{n+1} A_{n+1} + \eps_0 ^{-\beta}
\sum_{n\,\ge\,0} \left(
\frac{\eps_0}{\rho^{\beta+2}} \right) ^n
\| g \|_{L^\infty(Q_2)}.
\]
Assuming $\eps_0 < \rho^{\beta+2}<1$ and cancelling terms gives finally:
\[
A_0 \lesssim [S]_{\cC_\ell^{\alpha}(Q_{2})} + \| g \|_{L^\infty(Q_2)}
\]
which concludes the proof.
\end{proof}

\section{Global existence for the toy model}\label{sec:toy}

This section is devoted to the proof of Theorem~\ref{theo:main-holder}. Equation~\eqref{eq:main} is rewritten with the unknown function $g := f \mu^{-\frac12}$, where $\mu$ denotes the Gaussian $(2\pi)^{-d/2} e^{-|v|^2/2}$. This new function satisfies
\begin{equation}\label{eq:toyagain}
\partial_t g + v \cdot \nabla_x g = \fR[g] U[g]
\end{equation}
\begin{align*}
 \text{ with } \quad
\fR[g] &:= \int_v g \sqrt\mu \dd v \\
\text{ and } \quad U[g] &:=
\Delta_v g + \left(\frac{d}2 - \frac{|v|^2}{4}\right) g =
\frac1{\sqrt{\mu}} \nabla_v \left(\mu \nabla_v
\left(\frac{1}{\sqrt{\mu}} g\right)\right).
\end{align*}
After this rescaling the operator $U$ is symmetric in $L^2(\dd x \dd v)$, without weight. In contrast with~\eqref{eq:main}, this operator has no first order term in the velocity variable. But a (simpler) difficulty is created with the appearance of the unbounded zero order term $(\tfrac{d}2 - \tfrac{|v|^2}4)g$. We overcome it using that $g$ stays in between two Maxwellians (see Lemma~\ref{l:gaussian} below) and that, after rescaling, the H\"older estimate from~\cite{gimv,wz09} in Theorem~\ref{t:gimv} encodes decay in the $v$ variable (Proposition~\ref{prop:local}).

\subsection{Gaussian bounds}

We first explain how to propagate in time the Gaussian bounds satisfied by the initial data $g_{in} (x,v) = g(0,x,v)$.
\begin{lemm}[Gaussian bounds]\label{l:gaussian}
Consider a strong solution $g$ to~\eqref{eq:toyagain} in $L^\infty_t([0,T],H^2(\T^d\times\R^d))$ such that $C_1 \sqrt{\mu} \le g(0,x,v) \le C_2 \sqrt{\mu}$ in $\T^d \times \R^d$, then for almost all $t \in [0,T]$, we have
\[
\forall \, (x,v) \in \T^d \times \R^d, \quad C_1 \sqrt{\mu(v)} \le g(t,x,v) \le C_2 \sqrt{\mu(v)}.
\]
\end{lemm}


\begin{rema}
The notion of strong solutions used in the latter lemma and the rest of this section is standard: all the terms in the equation are defined almost everywhere and the equation is satisfied almost everywhere and the solution is continuous in time with value in $L^2$. We could have considered weaker notions of solutions but it was unnecessary since we are interested in constructing global
\emph{strong} solutions.
\end{rema}


\begin{proof}
For $a \in \R$, we write $a_+ = \max (0,a)$ and $a_- =
\max(0,-a)$. It is enough to prove that, for any $C>0$,
\[
\int_{\T^d\,\times\,\R^d} \left(g (t,x,v) - C \sqrt{\mu} (v)\right)^2_\pm \dd x \dd v
\]
decreases along time. It follows from the fact that $\tilde g = g - C\sqrt\mu$ satisfies
\[
\partial_t \tilde g + v \cdot \nabla_x \tilde g = \fR[g] U[\tilde g]
\]
and, multiplying the previous equation by $(\tilde g)_+ = \max (0,\tilde g)$ or $(\tilde g)_- = \max(0,-\tilde g)$, and integrating with respect to $x$ and $v$ yields for almost all $t \in [0,T]$,
\[
\int_{\T^d\,\times\,\R^d} \left(g (t,x,v) - C \sqrt{\mu} (v)\right)^2_\pm \dd x \dd v \le \int_{\T^d\,\times\,\R^d} \left(g_{in} (x,v) - C \sqrt{\mu} (v)\right)^2_\pm \dd x \dd v.
\]
In particular, for almost all $t \in [0,T]$,
\begin{align*}
\int_{\T^d\,\times\,\R^d} \left(g (t,x,v) - C_2 \sqrt{\mu} (v)\right)^2_+ \dd x \dd v \le \int_{\T^d\,\times\,\R^d} \left(g_{in} (x,v) - C_2 \sqrt{\mu} (v)\right)^2_+ \dd x \dd v = 0 \\
\int_{\T^d\,\times\,\R^d} \left(g (t,x,v) - C_1 \sqrt{\mu} (v)\right)^2_- \dd x \dd v \le \int_{\T^d\,\times\,\R^d} \left(g_{in} (x,v) - C_1 \sqrt{\mu} (v)\right)^2_- \dd x \dd v = 0.
\end{align*}
We conclude that for almost all $t \in [0,T]$, $C_1 \sqrt\mu (v) \le g (t,x,v) \le C_2 \sqrt\mu (v)$ in $\T^d \times \R^d$.
\end{proof}

\subsection{The local H\"older estimate for kinetic Fokker--Planck equations}\label{ss:holder}

We recall in this subsection the main result from~\cite{gimv,wz09} as stated in~\cite{gimv} and used in the proof of Theorem~\ref{theo:main-holder}. Consider the equation
\begin{equation}\label{e:lin-div}
\partial_t g + v \cdot \nabla_x g =
\nabla_v(A \nabla_v g) + S_0 \quad \text{in} \quad Q_{2r}(z_0)
\end{equation}
with $A=A(t,x,v)$ real symmetric matrix whose eigenvalues lie in $[\lambda,\Lambda]$ for $0 < \lambda\linebreak\le \Lambda$.

\begin{theo}[Local H\"older estimate]\label{t:gimv}
Given a cylinder $Q_{2r}(z_0) \subset \R \times \R^d \times \R^d$ and a function $S_0$ essentially bounded in $Q_{2r}(z_0)$, any weak solution $g$ of~\eqref{e:lin-div} in $Q_{2r}(z_0)$ such that $g \in L^\infty_t([0,T], L^2_{x,\,v} (\R^d \times \R^d)) \cap L^2_{t,\,x} ([0,T]
\times \R^d, H^1_v (\R^d))$ satisfies
\[
\| g\|_{\cC_\ell^{{\alpha_0}}\left(Q_{\frac{3r}2}(z_0)\right)} \le C r^{-{\alpha_0}} \| g \|_{L^2\left(Q_{2r}(z_0)\right)} + C r^{-2-{\alpha_0}}
\|S_0\|_{L^\infty\left(Q_{2r}(z_0)\right)}
\]
for some constants $\alpha_0 \in (0,1)$ and $C>0$ only depending on $d$ and $\lambda,\Lambda$.
\end{theo}


\subsection{The Schauder estimate for the toy model}\label{sec:scha-estim-toy}

The Schauder estimate follows from (1) the H\"older regularity established in \cite{gimv, wz09}, (2) the Schauder estimate from Section~\ref{sec:schauder} and (3) the Gaussian bounds from the previous subsection. Since the H\"older regularity only holds for positive times, we consider some time interval $[\tau,T]$ for some arbitrary but fixed $\tau>0$.

\begin{prop}[Higher order H\"older Estimates]\label{prop:local}
There exists ${\alpha_0} \in (0,1)$, only depending on the dimension $d$ and $C_1,C_2$, and $C>0$ only depending on $C_1,C_2,d, \tau,\delta$ such that for all solution $g \in \cC_\ell^{2+{\alpha_0}}$ to~\eqref{eq:toyagain} that satisfies $C_1 \sqrt \mu \le g \le C_2 \sqrt \mu$, for all $\tau \in (0,T)$ and $\delta \in (0,1)$, we have
\begin{equation}\label{eq:decay}
\| g \|_{\cC_\ell^{2+{\alpha_0}}\left(Q_r(z_0)\right)} \le C \mu^\delta (v_0).
\end{equation}
for any $Q_{2r} (z_0) \subset [\tau, T] \times \T^d \times \R^d$. In particular,
\begin{equation}\label{e:decay-grad}
\| \mu^{-\delta} \nabla_v g \|_{L^\infty\left([\tau,\,T] \,\times\,\R^d\,\times\,\R^d\right)} \le C.
\end{equation}
\end{prop}


\begin{proof}
We apply Theorem~\ref{t:gimv} with $\lambda = C_1$, $\Lambda = C_2$, $A = \fR[g] \Id$ and $S_0 = \fR[g](\frac{d}2-\frac{|v|^2}4)g$ and get
\[
\| g\|_{\cC_\ell^{{\alpha_0}}\left(Q_{\frac{3r}2}(z_0)\right)}
\le C r^{-{\alpha_0}} \| g \|_{L^2\left(Q_{2r}(z_0)\right)} + C r^{-2-{\alpha_0}}
\left\|\left(\frac{d}2 - \frac{|v|^2}4\right)g \right\|_{L^\infty\left(Q_{2r}(z_0)\right)}
\]
for some $\alpha_0 \in (0,1)$. Apply now Lemma~\ref{l:gaussian} to get for an arbitrary $\delta \in (0,1)$
\[
\| g\|_{\cC_\ell^{{\alpha_0}}\left(Q_{\frac{3r}2}(z_0)\right)} \le C_\delta \mu^{\delta}(v_0)
\]
for some constant $C_\delta$ depending only on $C_1, C_2, d,\tau, \delta$. Note that the role of the time $\tau>0$ is to ensure that the constant $C_\delta$ only depends on $\tau$ and not on $r$: it gives some ``room'' around a point $z_0 \in [\tau,T] \times \T^d \times \R^d$.

This implies in turn that for all $(t,x) \in [\tau,T] \times \T^d$,
\begin{align*}
\left|\fR[g](t,x) - \fR[g](s,y)\right| & \le \int_{\R^d} \left| g (t,x,v) - g (s,y,v)\right|\mu^{\frac12} (v) \dd v \\
& \le C_\delta \int_{\R^d} \left(|t-s|^{\frac{\alpha_0}2} + |x-y - (t-s)v|^{\frac{\alpha_0}3} \right) \mu^{\frac12+\delta} (v) \dd v \\
& \lesssim C_\delta \left(|t-s|^{\frac{\alpha_0}2} + |x-y|^{\frac{\alpha_0}3} + |t-s|^{\frac{\alpha_0}3}\right).
\end{align*}

This also implies for any $\delta' \in (0,\delta)$
\[
\left\|\left(\frac{d}2-\frac{|v|^2}4\right) g\right\| _{\cC_\ell^{{\alpha_0}}\left(Q_{\frac{3r}2}(z_0)\right)} \le C_\delta \mu^{\delta'}(v_0).
\]

This ensures the H\"older regularity of the coefficients and source term, and we are thus in a position to apply Theorem~\ref{thm:variable-loc} in the cylinders $Q_r (z_0)$ and deduce~\eqref{eq:decay}. To get~\eqref{e:decay-grad} from~\eqref{eq:decay} and the decay of Lemma~\ref{l:gaussian}, apply Lemma~\ref{l:interpol} on a cylinder $Q$ to obtain $\| \nabla_v g \|_{L^\infty(Q)} \le C \|g\|_{L^\infty(Q)} + C [g]_{\cC_\ell^{2+\alpha_0}(Q)}$ and argue as before. This achieves the proof of the Lemma~\ref{l:gaussian}.
\end{proof}


\subsection{Standard interpolation product inequality}\label{sec:stand-interp-prod}

We recall and prove an interpolation inequality tailored to our needs; it is folklore knowledge.

\begin{lemm}\label{lem:product}
Consider $k \ge d/2$, two functions $g_1, g_2 \in H^k(\T^d)$ and ${\bar m}, m \in \N^d$ such that $|{\bar m}| + |m| = k$ then
\[
\left\| \partial^{\bar m}_x g_1 \partial_x ^m g_2 \right\|_{L^2(\T^d)}
\lesssim_k \| g_1 \|_{L^\infty(\T^d)} \| g_2 \|_{H^k(\T^d)} +
\| g_1 \|_{H^k(\T^d)} \| g_2 \|_{L^\infty(\T^d)}.
\]
Moreover whenever ${\bar m} \not =0$, for any $\eps>0$ there is $C_\eps>0$ s.t.
\begin{equation}\label{eq:producteps}
\left\| \partial^{\bar m}_x g_1 \partial_x ^m g_2 \right\|_{L^2(\T^d)}
\le \eps \| g_1 \|_{L^\infty(\T^d)} \| g_2 \|_{H^k(\T^d)} + C_\eps \| g_1 \|_{H^k(\T^d)} \| g_2 \|_{L^\infty(\T^d)}.
\end{equation}
\end{lemm}


\begin{proof}
The first inequality is clear when $\bar m=0$. Use otherwise the \emph{Nash inequality}: given $h \in H^k(\T^d)$ with $k >d/2$ and any ${\bar m} \in \N^d$, with $0 < |{\bar m}| \le k$ then for $p := \frac{2k}{|{\bar m}|} \in [1,+\infty]$ one has
\[
\left\| \partial^{\bar m}_x h \right\|_{L^p(\T^d)} \lesssim \| h
\|_{L^\infty(\T^d)} ^{1-\frac{|{\bar m}|}k} \| h \|_{H^k(\T^d)} ^{\frac{|{\bar m}|}k}.
\]

Apply successively the H\"older inequality with $p := \frac{2k}{|{\bar m}|}$ and $q := \frac{2k}{|m|}$, and Nash inequality above:
\begin{align*}
\left\| \partial^{\bar m}_x g_1 \partial_x ^m g_2 \right\|_{L^2(\T^d)} & \le \left\| \partial^{\bar m}_x g_1 \right\|_{L^p(\T^d)} \left\|
\partial^{\bar m}_x g_2 \right\|_{L^q(\T^d)} \\
& \lesssim \| g_1\|_{L^\infty(\T^d)} ^{1-\frac{|{\bar m}|}{k}}
\| g_1 \|_{H^k(\T^d)}^{\frac{|{\bar m}|}k} \| g_2
\|_{L^\infty(\T^d)} ^{1-\frac{|m|}k} \| g_2 \|_{H^k(\T^d)} ^{\frac{|m|}k} \\
& \lesssim \left(\| g_1\|_{L^\infty(\T^d)} \| g_2
\|_{H^k(\T^d)}\right)^{1-\frac{|{\bar m}|}k}
\left(\| g_1 \|_{H^k(\T^d)} \| g_2\|_{L^\infty(\T^d)}
\right)^{\frac{|{\bar m}|}k}.
\end{align*}
Use finally $a^{1-\lambda} b^\lambda \le (1-\lambda)a + \lambda b$ to deduce the two claimed inequalities.
\end{proof}

\subsection{Local well-posedness in Sobolev spaces}\label{sec:local-wellp-sobol}

We prove that equation~\eqref{eq:main} is locally well-posed in $H^k (\T^d \times \R^d)$. It is convenient to work with the following norm,
\[
\| g \|_{H^{k,\,\bar k}_{x,\,v}\left(\T^d\,\times\,\R^d\right)} ^2 := \| g \|_{L^2\left(\T^d\,\times\,\R^d\right)} ^2 + \sum_{i=1,\,\dots,\,d} \left\| \partial_{x_i} ^k h\right\|_{L^2\left(\T^d\,\times\,\R^d\right)}^2 + \sum_{i=1,\,\dots,\,d} \left\| \partial_{v_i} ^{\bar k} h\right\|_{L^2\left(\T^d\,\times\,\R^d\right)}^2.
\]
Note that this norm agrees with the usual $H^k$ norm in $x$ and $v$ when $k = \bar k$: this is clear from the Fourier-transformed definition of these norms. However for $k \not = \bar k$ the space $H^{k,\,\bar k}_{x,\,v}$ includes but differs from $H^{\min(k,\,\bar k)}_{x,\,v}$. We use this intermediate definition of norm solely to shorten some calculations.


\begin{theo}[Local well-posedness in $H^k$]\label{t:loc-exis}
Consider
$
g_{in} \in H^{k,\,\bar k}_{x,\,v}(\T^d \times \R^d)
$
with $k,{\bar k}$ non-negative integers s.t. $2 \le{\bar k} \le k$ and $k > d/2$, and $C_1 \sqrt{\mu} \le g_{in} \le C_2 \sqrt{\mu}$ for $0 < C_1 < C_2$. Then there is $T>0$ depending only on $C_1, C_2$ and
\[
\left\|g_{in}\right\|_{H^{k,\,\bar k}_{x,\,v}\left(\T^d \times \R^d\right)}
\]
such that there exists a unique strong solution
\[
g \in C^1\left([0,T],H^{k,\,\bar k}_{x,\,v}\left(\T^d \times \R^d\right)\right)
\]
to~\eqref{eq:toyagain} with initial data $g_{in}$, which furthermore satisfies $C_1 \sqrt{\mu} \le g (t,\cdot,\cdot) \le C_2 \sqrt{\mu}$ for almost all $t \in [0,T]$.
\end{theo}


\begin{proof}
It is standard calculations that the two a priori estimates~\eqref{eq:gronwall} and~\eqref{eq:difference} below imply the existence of solutions constructed through the iterative scheme
\begin{equation}\label{eq:iterated}
\begin{cases}
(\partial_t + v \cdot \nabla_x) g_{n+1} &\; =\; \fR [g_n] U[g_{n+1}],\\
 g_{n+1}(t=0,\cdot,\cdot)&\; =\; g_{in} \\
\text{Initialization: } g_0(t,x,v) & :=\; g_{in}(x,v)  \forall \, t \ge 0.
\end{cases}
\end{equation}

We now focus on establishing the key a priori estimate. Consider a solution $g_{n+1} \in C^1([0,T],H^{k,\,\bar k}_{x,\,v}(\T^d \times \R^d))$ to~\eqref{eq:iterated} and compute successively:
\begin{equation}\label{eq:evol-2}
L^2 \text{ estimate: } \qquad \dt \frac{1}2 \int_{\T^d\, \times\,\R^d} | g_{n+1}|^2\dd x \dd v \le - C_1
\int_{\T^d\,\times\,\R^d} |h_{n+1}|^2 \dd x \dd v
\end{equation}
where we denote $h_{n+1}:=\mu^{1/2} \nabla_v (\mu^{-1/2} g_{n+1})$. Regarding the $v$-derivatives, for any integer ${\bar k} \ge 1$,
\begin{align*}
\dt \frac12 \int_{\T^d\,\times\,\R^d} \left|\partial^{{\bar k}}_{v_i} g_{n+1}\right|^2 \dd x \dd v = &- {\bar k} \int_{\T^d\, \times\,\R^d} \left(\partial^{{\bar k}-1} _{v_i} \partial_{x_i} g_{n+1}\right) \partial^{{\bar k}}_{v_i} g_{n+1} \dd x \dd v \\
&- \int_{\T^d\,\times\,\R^d} \fR[g_n] \left| \nabla_v \left(\frac{\partial^{{\bar k}} _{v_i} g_{n+1}}{\sqrt \mu} \right) \right|^2 \mu \dd x \dd v\\
& + \frac14 \binom{{\bar k}}1 \int_{\T^d\,\times\,\R^d} \fR[g_n] \left| \partial^{{\bar k} - 1} _{v_i} g_{n+1} \right|^2 \dd x \dd v \\
& + \frac12 \binom{{\bar k}}2 \int_{\T^d\,\times\,\R^d} \fR[g_n] \left| \partial^{{\bar k}-1} _{v_i} g_{n+1} \right|^2 \dd x \dd v.
\end{align*}
In the right hand side, the first term corresponds to the transport $v \cdot \nabla_x$, the second one to the operator $U$ since $\fR[g_n]$ does not depend on $v$, the third term appears when one $v$-derivative applies to $|v|^2$ and the others apply to $g_{n+1}$ in the product $|v|^2 g_{n+1}$ appearing in $U[g_{n+1}]$, the fourth term appears after deriving $|v|^2$ twice. Notice that integrations by parts are used either to further differentiate $|v|^2$ or to make appear $|\partial_{v_i}^{{\bar k}-1}g_{n+1}|^2$. Discarding the negative term and using the fact that $\fR[g_n] \le C_2$ thanks to Lemma~\ref{l:gaussian}, we get after summing over $i =1,\,\dots,\,d$ and combining with equation~\eqref{eq:evol-2}
\begin{multline}\label{eq:evol-v}
\text{Estimate of $v$-derivatives}:\\
\dt \frac12 \left\|g_{n+1}(t) \right\|^2_{H^{0,\,\bar k}_{x,v} \left(\T^d\,\times\,\R^d\right)}  \lesssim_{{\bar k},\,C_2}
\left\| g_{n+1} \right\|_{H^{\bar k,\,\bar k}_{x,v}\left(\T^d\,\times\,\R^d\right)} ^2.
\end{multline}
Regarding the $x$-derivatives, we write the equation on $g_{n+1}$ as
\[
(\partial_t + v \cdot \nabla_x) g_{n+1} = \fR [g_n] \mu^{-\frac12}
\nabla_v \cdot \left(\mu^{\frac12} \nabla_v h_{n+1} \right) \quad
\text{ with } \quad h_{n+1}:= \mu^{\frac12} \,\nabla_v \left(\mu^{-\frac12} g_{n+1}\right).
\]
Since $x$-derivatives commute with the operators $v\cdot
\nabla_x$ and $U$, we have for all\linebreak$i=1,\,\dots,\,d$,
\begin{align*}
\dt \frac12 \int_{\T^d \times \R^d} \left|\partial^{k}_{x_i} g_{n+1}\right|^2 \dd x \dd v &=  - \int_{\T^d\, \times\,\R^d} \fR[g_n] \left| \partial_{x_i}^k h_{n+1}\right|^2 \dd x \dd v\\
& -\!\! \sum_{0\,\le\,q\,<\,k} \binom{k}{q} \int_{\T^d\,\times \,\R^d} \left(\partial^{k - q}_{x_i} \fR[g_n]\right) \partial^{q}_{x_i} h_{n+1} \cdot \partial^{k}_{x_i} h_{n+1} \dd x \dd v.
\end{align*}
Use now $\fR[g_n] \ge C_1$ from Lemma~\ref{l:gaussian}:
\begin{align*}
\dt \frac12 \int_{\T^d\,\times\,\R^d} \left|\partial^{k}_{x_i} g_{n+1}\right|^2
\dd x \dd v &\lesssim_k  - C_1 \int_{\T^d\,\times\,\R^d} \left| \partial_{x_i}^{k}h_{n+1}\right|^2 \dd x \dd v \\
& +\!\! \sum_{0\,\le\,q\,<\,k} \int_{\T^d\,\times\,\R^d} \left|\partial^{k - q} _{x_i}
\fR[g_n]\right|\! \cdot\!\left|\partial_{x_i}^{q} h_{n+1}\right|\!
\cdot\!\left|\partial_{x_i}^{k} h_{n+1}\right| \dd x \dd v.
\end{align*}
Observe that, given $0 \le q < k$, the index $p:=k-q \not =0$ and the inequality~\eqref{eq:producteps} in Lemma~\ref{lem:product} can be applied (we use again below the upper bound $\fR[g_n] \le C_2$):
\begin{multline*}
\int_{\T^d \times \R^d} \left|\partial^{k-q} _{x_i}
\fR[g_n]\right|
\cdot \left|\partial_{x_i}^{q} h_{n+1}\right| \cdot \left|\partial_{x_i}^{k} h_{n+1}\right| \dd x \dd v \\
\begin{aligned}
& = \int_{\R^d} \left(\int_{\T^d} \left|\partial^{p} _{x_i}
\fR[g_n]\right|
\cdot \left|\partial_{x_i}^{q} h_{n+1}\right|
\cdot \left|\partial_{x_i}^{k} h_{n+1}\right| \dd x \right) \dd v \\
& \le \int_{\R^d} \left(\left\| \partial^{p} _{x_i} \fR[g_n](t,\cdot) \partial_{x_i}^{q} h_{n+1}(t,\cdot,v) \right\|_{L^2_x(\T^d)}
\left\|\partial_{x_i}^{k} h_{n+1}(t,\cdot,v) \right\|_{L^2_x(\T^d)} \right) \dd v \\
&\le \eps \int_{\R^d} \left(\left\| \fR[g_n](t,\cdot)\right\|_{L^\infty_x(\T^d)} \left\| h_{n+1}(t,\cdot,v) \right\|_{H^k_x(\T^d)} ^2\right) \dd v \\
& + C_\eps \int_{\R^d} \left(\left\| \fR[g_n](t,\cdot)\right\|_{H^k _x(\T^d)} \left\| h_{n+1}(t,\cdot,v) \right\|_{L^\infty_x(\T^d)} \left\| h_{n+1}(t,\cdot,v) \right\|_{H^k_x(\T^d)} \right) \dd v\\
&\le \eps \left(C_2 + 1\right)
\left\| h_{n+1}(t,\cdot,\cdot) \right\|_{H^{k,0}_{x,v}\left(\T^d\,\times\,\R^d\right)} ^2 \\
&+ C_\eps '\left\| \fR[g_n](t,\cdot)\right\|_{H^k _x(\T^d)} ^2
\int_{\R^d} \left\| h_{n+1}(t,\cdot,v) \right\|_{L^\infty_x(\T^d)} ^2 \dd v
\end{aligned}
\end{multline*}
for any $\eps >0$ and some corresponding constant $C_\eps,C_\eps'>0$. Use then
\[
\left\| \fR[g_n](t,\cdot)\right\|_{H^k _x(\T^d)} \le \left\| g_n(t,\cdot,\cdot)\right\|_{H^{k,\,0}_{x,\,v}\left(\T^d\times \R^d\right)}
\]
and equation~\eqref{eq:evol-2} to get
\begin{multline*}
\dt \frac12 \|g_{n+1}\|^2_{H^{k,\,0}_{x,\,v}\left(\T^d \times \R^d\right)}\\
\begin{aligned}
&\lesssim_{k,\,C_2} - C_1 \left\| h_{n+1}(t,\cdot,\cdot)\right\|_{H^{k,\,0}_{x,\,v}\left(\T^d \times
\R^d\right)} ^2 + \eps \left\| h_{n+1}(t,\cdot,\cdot)\right\|_{H^{k,\,0}_{x,\,v}\left(\T^d \times
\R^d\right)} ^2 \\
&\quad+ C'_\eps \left(\int_{\R^d} \left\| h_{n+1}(t,\cdot,v) \right\|_{L^\infty_x(\T^d)} ^2 \dd v \right) \left\| g_n(t,\cdot,\cdot)\right\|_{H^{k,\,0}_{x,\,v}\left(\T^d \times \R^d\right)} ^2.
\end{aligned}
\end{multline*}
Finally choose $\eps$ small enough (in terms of absolute constants, independently of the solution) to get
\begin{multline}\label{eq:evol-x}
\dt \frac12 \|g_{n+1}\|^2_{H^{k,\,0}_{x,v}\left(\T^d\,\times\,\R^d\right)} \\
 \lesssim_{k,\,C_1,\,C_2} \left(\int_{\R^d} \left\| h_{n+1}(t,\cdot,v) \right\|_{L^\infty_x(\T^d)}^2 \dd v \right)
\left\| g_n(t,\cdot,\cdot) \right\|_{H^{k,\,0}_{x,\,v}\left(\T^d\,\times\,\R^d\right)}^2.
\end{multline}

The combination of equations~\eqref{eq:evol-v} and~\eqref{eq:evol-x} yields
\begin{multline}\label{eq:evol-global}
\dt \frac12 \left\| g_{n+1}(t,\cdot,\cdot)\right\|_{H^{k,\,\bar k}_{x,\,v}\left(\T^d\,\times\,\R^d\right)}^2 \lesssim_{k,C_1,\,C_2}
\left\| g_{n+1} (t,\cdot,\cdot) \right\|_{H^{\bar k,\,\bar k}_{x,v}\left(\T^d\,\times\,\R^d\right)}^2 \\
+ \left(\int_{\R^d} \left\| h_{n+1}(t,\cdot,v)
\right\|_{L^\infty_x(\T^d)}^2 \dd v \right) \left\| g_n(t,\cdot,\cdot) \right\|_{H^{k,\,0}_{x,\,v}\left(\T^d\,\times
\,\R^d\right)}^2.
\end{multline}

Now observe that $h_{n+1} = \nabla_v g_{n+1} + \frac{v}{2} g_{n+1}$ and use $g_{n+1} \le C_2 \sqrt\mu$ and $k>d/2$ and ${\bar k} \ge 1$ and Sobolev embedding in $\T^d$ to get
\begin{align*}
\int_{\R^d} \left\| h_{n+1}(t,\cdot,v) \right\|_{L^\infty_x(\T^d)} ^2 \dd v & \lesssim C_2^2+ \int_{\R^d} \| \nabla_v g_{n+1} \|_{L^\infty_x(\T^d)} ^2 \dd v \\
& \lesssim C_2^2 + \| g_{n+1} \|_{H^{k,\,\bar k}_{x,\,v} \left(\T^d\,\times\,\R^d\right)}^2
\end{align*}
and conclude finally that
\begin{multline}\label{eq:gronwall}
\dt \frac12 \left\| g_{n+1}(t,\cdot,\cdot)\right\|_{H^{k,\,\bar k}_{x,\,v}\left(\T^d\,\times\,\R^d\right)}^2 \\
\lesssim_{k,\,C_1,\,C_2}
\left(\left\| g_n(t,\cdot,\cdot) \right\|_{H^{k,\,\bar k}_{x,\,v}\left(\T^d\,\times\,\R^d\right)}^2 + 1 \right)
\left\| g_{n+1}(t,\cdot,\cdot) \right\|_{H^{k,\,\bar k}_{x,\,v}\left(\T^d\,\times\,\R^d\right)}^2
\end{multline}
which is the first main a priori estimate, that shows that the $H^{k,\,\bar k}_{x,\,v}(\T^d\times \R^d)$ norm remains finite on a short time interval (whose length depends on the size of the initial data) thanks to Gronwall's lemma.
%\medskip

Regarding uniqueness, consider the difference of two solutions $g,{\bar g} \in H^{k,\,\bar k}_{x,\,v}(\T^d \times \R^d)$ that satisfies
\[
\partial_t(g-{\bar g}) + v \cdot \nabla_x (g-{\bar g}) =
\fR[g-{\bar g}]U[g]+ \fR[{\bar g}] U[g-{\bar g}],
\]
and perform similar energy estimates to get
\begin{multline}\label{eq:difference}
\dt \frac12 \left\| (g-{\bar g})(t,\cdot,\cdot)
\right\|_{H^k_{x}H^{\bar k}_v\left(\T^d\,\times\,\R^d\right)}^2
\lesssim_{k,\,C_1,\,C_2} \left\| (g-{\bar g})(t,\cdot,\cdot)\right\|_{H^{k,\,\bar k}_{x,v}\left(\T^d\,\times \,\R^d\right)}^2 \\
+ \left(\left\| g(t,\cdot,\cdot)
\right\|_{H^{k,\,\bar k}_{x,\,v}\left(\T^d\,\times\,\R^d\right)}^2 + \left\| {\bar g}(t,\cdot,\cdot)
\right\|_{H^{k,\,\bar k}_{x,\,v}\left(\T^d\,\times\,\R^d\right)} ^2 \right)
\left\| (g-{\bar g})(t,\cdot,\cdot)
\right\|_{H^{k,\,\bar k}_{x,\,v}\left(\T^d\,\times\,\R^d\right)}^2
\end{multline}
which implies uniqueness in the space $H^{k,\,\bar k}_{x,\,v}(\T^d\,\times\,\R^d)$ and concludes the proof of Theorem~\ref{t:loc-exis}.
\end{proof}

\subsection{From local-in-time to global-in-time}\label{sec:from-local-time}

To continue the solutions for all times it is enough to prove that the $H^{k,\,\bar k}_{x,\,v}(\T^d \times \R^d)$ norm remains finite over arbitrarily long times. Thanks to the local-in-time existence result in Theorem~\ref{t:loc-exis}, we know that solutions exist at least on a time interval $[0,2 \tau]$ for some $\tau>0$. We now show that the Sobolev norm cannot explode after time $\tau$ by using Proposition~\ref{prop:local}.

Consider again the a priori estimate~\eqref{eq:evol-global},
\begin{multline*}
\dt \frac12 \left\| g_{n+1}(t,\cdot,\cdot) \right\|_{H^{k,\,\bar k}_{x,v}\left(\T^d\,\times\,\R^d\right)}^2  \lesssim_{k,\,C_1,\,C_2} \left\| g_{n+1}(t,\cdot,\cdot)
\right\|_{H^{\bar k,\,\bar k}_{x,\,v}\left(\T^d\,\times\, \R^d\right)}^2 \\
+ \left(\int_{\R^d} \left\| h_{n+1}(t,\cdot,v)
\right\|_{L^\infty_x(\T^d)}^2 \dd v \right)
\left\| g_n(t,\cdot,\cdot) \right\|_{H^{k,\,0}_{x,\,v}\left(\T^d\,\times\,\R^d\right)}^2.
\end{multline*}

Recall that $h_{n+1} = \nabla_v g_{n+1} + \frac{v}{2} g_{n+1}$ and use Lemma~\ref{l:gaussian}:
\[
\left(\int_{\R^d} \left\| h_{n+1}(t,\cdot,v)
\right\|_{L^\infty_x(\T^d)}^2 \dd v \right)
\lesssim \left(\int_{\R^d} \left\| \nabla_v g_{n+1}(t,\cdot,v)
\right\|_{L^\infty_x(\T^d)}^2 \dd v \right) + C_2.
\]
Apply finally Proposition~\ref{prop:local}:
\[
\left(\int_{\R^d} \left\|\nabla_v g_{n+1}(t,\cdot,v) \right\|_{L^\infty_x(\T^d)}^2 \dd v
\right) \lesssim \left\| \mu^{-\delta} \nabla_v g_{n+1} \right\|^2_{L^\infty\left([\tau,\,T]\,\times\,
\T^d\,\times\,\R^d\right)} \lesssim 1.
\]

It shows for $t \ge \tau$, i.e. after some arbitrarily small time $\tau>0$:
\begin{multline*}
\dt \frac12 \left\| g_{n+1}(t,\cdot,\cdot) \right\|_{H^{k,\,\bar k}_{x,\,v}\left(\T^d\,\times\,\R^d\right)}^2\\
\lesssim_{k,\,C_1,\,C_2,\,\tau} \left\| g_{n+1}(t,\cdot,\cdot)
\right\|_{H^{k,\,\bar k}_{x,v}\left(\T^d\,\times\,\R^d\right)}^2 +
\left\| g_n(t,\cdot,\cdot) \right\|_{H^{k,\,\bar k}_{x,\,v}\left(\T^d\,\times\,\R^d\right)}^2.
\end{multline*}
This proves that the norm remains finite by the Gronwall lemma and concludes the proof.



\begin{thebibliography}{99}
\bibitem{baudoin2017}Baudoin, Fabrice Bakry–Émery meet Villani, J. Funct. Anal., Volume 273 (2017) no. 7, pp. 2275-2291 | DOI | MR | Zbl
\bibitem{bb}Bramanti, Marco; Brandolini, Luca Schauder estimates for parabolic nondivergence operators of Hörmander type, J. Differ. Equations, Volume 234 (2007) no. 1, pp. 177-245 | DOI | MR | Zbl
\bibitem{MR763762}Beale, J. Thomas; Kato, Tosio; Majda, Andrew J. Remarks on the breakdown of smooth solutions for the $3$-D Euler equations, Commun. Math. Phys., Volume 94 (1984) no. 1, pp. 61-66 | MR | Zbl | DOI
\bibitem{boltzmann1872weitere}Boltzmann, Ludwig Weitere Studien über das Wärmegleichgewicht unter Gasmolekülen, Sitzungs. Further Studies on the Thermal Equilibrium of Gas Molecules, Kinetic Theory, Volume 2 (1872), pp. 88-174 | Zbl
\bibitem{MR889476}Bakry, Dominique; Émery, Michel Diffusions hypercontractives, Séminaire de probabilités, XIX, 1983/84 (Lecture Notes in Mathematics), Volume 1123, Springer, 1985, pp. 177-206 | DOI | MR | Zbl | Numdam
\bibitem{chavanis2008nonlinear}Chavanis, Pierre-Henri Nonlinear mean field Fokker–Planck equations. Application to the chemotaxis of biological populations, Eur. Phys. J. B, Condens. Matter Complex Syst., Volume 62 (2008) no. 2, pp. 179-208 | Zbl | DOI
\bibitem{dfp}Di Francesco, Marco; Polidoro, Sergio Schauder estimates, Harnack inequality and Gaussian lower bound for Kolmogorov-type operators in non-divergence form, Adv. Differ. Equ., Volume 11 (2006) no. 11, pp. 1261-1320 | MR | Zbl
\bibitem{gimv}Golse, François; Imbert, Cyril; Mouhot, Clément; Vasseur, Alexis F. Harnack inequality for kinetic Fokker–Planck equations with rough coefficients and application to the Landau equation, Ann. Sc. Norm. Super. Pisa, Cl. Sci., Volume 19 (2019) no. 1, pp. 253-295 | MR | Zbl
\bibitem{guillin2012}Guillin, Arnaud; Wang, Feng-Yu Degenerate Fokker–Planck equations: Bismut formula, gradient estimate and Harnack inequality, J. Differ. Equations, Volume 253 (2012) no. 1, pp. 20-40 | DOI | MR | Zbl
\bibitem{hs}Henderson, Christopher; Snelson, Stanley ${C}^\infty $ smoothing for weak solutions of the inhomogeneous Landau equation, Arch. Ration. Mech. Anal., Volume 236 (2020) no. 1, pp. 113-143 | Zbl | DOI | MR
\bibitem{bootstrap}Imbert, Cyril; Silvestre, Luis Global regularity estimates for the Boltzmann equation without cut-off (2019) (https://arxiv.org/abs/1909.12729)
\bibitem{is}Imbert, Cyril; Silvestre, Luis The weak Harnack inequality for the Boltzmann equation without cut-off, J. Eur. Math. Soc., Volume 22 (2020) no. 2, pp. 507-592 | Zbl | MR | DOI
\bibitem{imbert2018schauder}Imbert, Cyril; Silvestre, Luis The Schauder estimate for kinetic integral equations, Anal. PDE, Volume 14 (2021) no. 1, pp. 171-204 | DOI | MR
\bibitem{MR0084985}Kac, Mark Foundations of kinetic theory, Proceedings of the Third Berkeley Symposium on Mathematical Statistics and Probability, 1954–1955, vol. III, University of California Press, Berkeley and Los Angeles (1956), pp. 171-197 | MR | Zbl
\bibitem{MR2252155}Kiessling, Michael; Lancellotti, Carlo The linear Fokker–Planck equation for the Ornstein–Uhlenbeck process as an (almost) nonlinear kinetic equation for an isolated ${N}$-particle system, J. Stat. Phys., Volume 123 (2006) no. 3, pp. 525-546 | DOI | MR | Zbl
\bibitem{kolm}Kolmogoroff, Andreĭ N. Zufällige Bewegungen (zur Theorie der Brownschen Bewegung), Ann. Math., Volume 35 (1934) no. 1, pp. 116-117 | MR | Zbl | DOI
\bibitem{krylov1996}Krylov, Nicolaĭ V. Lectures on elliptic and parabolic equations in Hölder spaces, Graduate Studies in Mathematics, 12, American Mathematical Society, 1996 | DOI | MR | Zbl
\bibitem{landau1936}Landau, Lev D. The transport equation in the case of Coulomb interactions – Die kinetische Gleichung für den Fall Coulombscher Wechselwirkung, Phys. Z. Sowjetunion, Volume 10 (1936), pp. 154-164 | Zbl
\bibitem{lunardi1997}Lunardi, Alessandra Schauder estimates for a class of degenerate elliptic and parabolic operators with unbounded coefficients in ${\bf {R}}^n$, Ann. Sc. Norm. Super. Pisa, Cl. Sci., Volume 24 (1997) no. 1, pp. 133-164 | MR | Zbl
\bibitem{manfredini}Manfredini, Maria The Dirichlet problem for a class of ultraparabolic equations, Adv. Differ. Equ., Volume 2 (1997) no. 5, pp. 831-866 | MR | Zbl
\bibitem{clerk1867dynamical}Maxwell, James Clerk On the dynamical theory of gases, Philosophical Transactions of the Royal Society of London Series I, Volume 157 (1867), pp. 49-88
\bibitem{polidoro1994}Polidoro, Sergio On a class of ultraparabolic operators of Kolmogorov–Fokker–Planck type, Matematiche, Volume 49 (1994) no. 1, pp. 53-105 | MR | Zbl
\bibitem{radkevich2008}Radkevich, Evgeniĭ V. Equations with nonnegative characteristic form. II, Sovrem. Mat. Prilozh. (2008) no. 56, pp. 3-147 (Differentsial' nye Uravneniya s Chastnymi Proizvodnymi) | DOI | MR | Zbl
\bibitem{safonov1984}Safonov, Mikhail V. The classical solution of the elliptic Bellman equation, Dokl. Akad. Nauk SSSR, Volume 278 (1984) no. 4, pp. 810-813 | MR
\bibitem{MR3551261}Silvestre, Luis A new regularization mechanism for the Boltzmann equation without cut-off, Commun. Math. Phys., Volume 348 (2016) no. 1, pp. 69-100 | DOI | MR | Zbl
\bibitem{villani2009memoir}Villani, Cédric Hypocoercivity, Memoirs of the American Mathematical Society, 950, American Mathematical Society, 2009 | DOI | MR | Zbl
\bibitem{wz09}Wang, WenDong; Zhang, LiQun The ${C}^\alpha $ regularity of a class of non-homogeneous ultraparabolic equations, Sci. China, Ser. A, Volume 52 (2009) no. 8, pp. 1589-1606 | Zbl | DOI | MR
\end{thebibliography}
\end{document}
